\documentclass[11pt]{article}
\usepackage{paper}

\newcommand{\repourl}{\textless\textit{repository URL to be inserted}\textgreater}

\title{Integral sums of reciprocals of distinct squarefree semiprimes:\\
forty-seven terms are necessary, and the\\ extremal representations of $1$}
% ------------------------------------------------------------------
% AUTHORS: intentionally left blank -- to be filled in by the authors.
\author{}
% ------------------------------------------------------------------
\date{}

\begin{document}
\maketitle

\begin{abstract}
Let $\Pset=\{pq: p<q \text{ primes}\}$ be the set of squarefree semiprimes.
We prove that if $T\subset\Pset$ is a nonempty finite set such that $\sum_{n\in T}1/n$ is an integer, then $\card{T}\ge 47$.
Since $47$-term representations of $1$ are known, $47$ is the exact minimum; this confirms a conjecture of Watanabe.
We also classify the extremal case: there are exactly $23$ sets $T\subset\Pset$ with $\card{T}=47$ and $\sum_{n\in T}1/n=1$.
Seventeen of them involve only primes at most $73$, and each of the other six involves exactly one prime larger than $73$, namely $269$, $311$, $421$, $433$, $503$ or $3779$.
The proofs are computer-assisted.
A $p$-adic argument turns integrality into a congruence condition at every vertex of a graph on the primes.
A decomposition into a ``core'' of small primes and a ``large-prime part'', together with value, parity and star-size lemmas and an explicit divisor equation for a single edge between two large primes, reduces the problem to a finite exhaustive computation, even though arbitrarily large prime factors are allowed.
The enumeration of the cores uses a new pruning bound, a Lagrangian relaxation of the value condition combined with a knapsack over residues, which reduces the size of the search tree by a factor of about $3000$ for $47$ terms.
The computation uses exact integer arithmetic and directed rounding only.
With the same method we also determined all $48$-term representations of $1$: there are exactly $620$ of them.
We describe independent checks: a second implementation of the completion step, two implementations of the core enumeration with identical outputs, brute-force validation on reduced instances, and exact verification of all solutions.
The analytic reduction and the $23$ solutions are verified in Lean~4 with Mathlib. The core enumeration is written as a Lean function, and Lean proves that this search is exhaustive: its output contains the core of every solution.
The remaining completion step enters the Lean proofs as explicitly stated hypotheses about the cores in that output; these hypotheses are established by computation and are not verified in Lean.
\end{abstract}

\tableofcontents

% =====================================================================
\section{Introduction}\label{sec:intro}
% =====================================================================

An \emph{Egyptian fraction} is a finite sum of distinct unit fractions. Many classical questions ask which rationals are sums of reciprocals of distinct integers from a prescribed set; see \cite{Graham1964,ErdosGraham1980,Graham2013} and \cite[Section~D11]{Guy2004}.
This paper studies the set
\[
\Pset=\{pq : p<q \text{ primes}\}=\{6,10,14,15,21,22,26,33,34,35,38,39,\dots\}
\]
of \emph{squarefree semiprimes}, i.e.\ the integers with exactly two distinct prime factors, each to the first power.
For a finite set $T\subset\Pset$ write
\[
\Rec{T}=\sum_{n\in T}\frac1n .
\]

\paragraph{Background.}
According to the abstract of Watanabe's preprint \cite{Watanabe2020}, in 1978 Johnson found a representation of $1$ as a sum of reciprocals of $48$ distinct squarefree semiprimes \cite{Johnson1978}. No shorter representation was known until Watanabe \cite{Watanabe2020} found $17$ representations with $47$ terms.
Watanabe found no representation with fewer than $47$ terms and conjectured that $47$ is the minimum.
His examples were obtained by computer searches in which the prime factors were bounded.\footnote{The search space described in \cite{Watanabe2020} has dimension $325=\binom{26}{2}$, the number of squarefree semiprimes whose prime factors are both at most $101$. We could not consult the full text while preparing this manuscript, so this detail should be checked against \cite{Watanabe2020}.}
A search with bounded prime factors cannot establish minimality, because arbitrarily large primes are allowed in principle; \cref{sec:stars} shows that primes as large as $3779$ really do occur.

In a different direction, Butler, Erd\H{o}s and Graham \cite{ButlerErdosGraham2015} proved that every positive integer is a sum of reciprocals of distinct integers that each have exactly three distinct prime factors, and they conjectured the corresponding statement for two prime factors. A recent preprint of Li \cite{Li2026} announces a proof of that conjecture.
Egyptian fractions whose denominators are supported on a fixed finite set of primes are studied computationally in \cite{Czenky2025}.
None of these works addresses the minimal number of terms for the set $\Pset$.
To the best of our knowledge, no proof that $47$ is minimal, and no complete list of the $47$-term representations of $1$ in which every prime is allowed, has appeared before.

\paragraph{Results.}

\begin{maintheorem}[Minimum number of terms]\label{thm:A}
Let $T\subset\Pset$ be a nonempty finite set such that $\Rec{T}\in\Z$. Then $\card{T}\ge 47$.
The bound is sharp: there are sets $T\subset\Pset$ with $\card{T}=47$ and $\Rec{T}=1$.
\end{maintheorem}

\begin{maintheorem}[The extremal representations]\label{thm:B}
There are exactly $23$ sets $T\subset\Pset$ with $\card{T}=47$ and $\Rec{T}=1$. They are listed in \cref{app:solutions}.
Equivalently, there are exactly $23$ sets $T\subset\Pset$ with $\card{T}=47$ and $\Rec{T}\in\Z$.
Of these $23$ sets:
\begin{enumerate}[label=\textup{(\roman*)}]
\item $17$ involve only primes at most $73$ (in fact at most $71$);
\item each of the other $6$ involves exactly one prime $P>73$, where $P$ is one of $269$, $311$, $421$, $433$, $503$ and $3779$. This prime occurs in exactly three elements $aP,bP,cP$, where $a,b,c\le 73$ are primes with $P\mid ab+ac+bc$ (see \cref{tab:bigprimes}).
\end{enumerate}
\end{maintheorem}

In particular, the $17$ solutions of (i) are exactly the solutions all of whose prime factors are at most $101$.
Our $17$ solutions of type (i) have the same count as Watanabe's list of $17$ examples. We have not compared the two lists term by term.

The method is not limited to $47$ terms. Run with budget $48$, the same search and completion determine all representations of $1$ by $48$ distinct squarefree semiprimes.

\begin{maintheorem}[Forty-eight terms]\label{thm:C}
There are exactly $620$ sets $T\subset\Pset$ with $\card{T}=48$ and $\Rec{T}=1$ (equivalently, $\Rec{T}\in\Z$). They are listed in the file \file{solutions48.txt} of the repository. Of these $620$ sets:
\begin{enumerate}[label=\textup{(\roman*)}]
\item $397$ involve only primes at most $73$;
\item $187$ involve exactly one prime $P>73$, all of whose neighbours are at most $73$;
\item $18$ involve exactly two primes $P,P'>73$, neither of which divides an element of $T$ together with the other;
\item $18$ involve exactly two primes $P_1,P_2>73$ and contain the element $P_1P_2$.
\end{enumerate}
\end{maintheorem}

\cref{thm:C} rests on one further computational ingredient, the treatment of shapes with two large--large elements (\cref{sec:proof48}), which is not needed for \cref{thm:A,thm:B}. The largest prime that occurs in a $48$-term representation is $1\,210\,883$ (\cref{sec:proof48}).

\paragraph{Nature of the proof.}
All three theorems are \emph{computer-assisted}.
\cref{sec:prelim,sec:graph,sec:core,sec:stars} give a complete mathematical reduction of the statements to finite computations.
We state the computations for \cref{thm:A,thm:B} precisely as \cref{claim:46,claim:47}.
The computation has two steps: the enumeration of all possible ``cores'' (Step~1) and their completion (Step~2).
Step~1 is carried out by a depth-first search with a new pruning bound, the \emph{loss bound} (\cref{sec:lossbound}); it is written as a Lean function, and Lean proves that the search is exhaustive (\cref{sec:lean}).
Both steps were executed by the C++ program \file{semiprime48.cpp}, which uses exact integer arithmetic and directed rounding only (\cref{sec:proof46,sec:proof47,sec:proof48,sec:repro}).
An earlier implementation without the loss bound (\file{semiprime_reciprocals.cpp} and \file{semiprime47.cpp}) produced identical results for \cref{thm:A,thm:B}.
\cref{sec:verification} describes the independent checks, and \cref{sec:lean} describes the Lean~4 formalization.
\Cref{fig:status} summarizes the logical status of each ingredient.

\begin{figure}[t]
\centering
\begin{tikzpicture}[
  box/.style={draw,rounded corners=2pt,align=center,font=\small,inner sep=5pt,text width=5.3cm,minimum height=1.25cm},
  lab/.style={font=\scriptsize\itshape,align=center},
  arr/.style={-{Stealth[length=2.2mm]},thick}]
\node[box] (red) {Analytic reduction\\ (\cref{sec:prelim,sec:graph,sec:core,sec:stars})};
\node[box,right=1.6cm of red] (step1) {Step 1: the core search\\ with the loss bound (\cref{sec:step1})};
\node[box,below=1.5cm of step1] (claim) {Step 2 on the cores found\\ \cref{claim:46,claim:47} (\lean{Search46'}, \lean{Search47'})};
\node[box,below=1.5cm of red] (sols) {The $23$ explicit solutions\\ (\cref{app:solutions})};
\node[lab,below=1pt of sols] (solslab) {verified exactly: C++, Python, Lean (unconditional)};
\node[lab,below=1pt of claim] (claimlab) {computed in C++ and in Python; \emph{not} verified in Lean};
\node[box,below=1.4cm of $(solslab.south)!0.5!(claimlab.south)$] (thm) {\cref{thm:A,thm:B}};
\draw[arr] (step1) -- (claim);
\draw[arr] (claimlab.south) -- ++(0,-0.3) -| ([xshift=0.9cm]thm.north);
\draw[arr] (solslab.south) -- ++(0,-0.3) -| ([xshift=-0.9cm]thm.north);
\draw[arr] (red.west) -- ++(-0.35,0) |- (thm.west);
\node[lab,above=1pt of red] {verified in Lean 4 (unconditional)};
\node[lab,above=1pt of step1] {search defined in Lean; exhaustiveness\\ proved in Lean (unconditional);\\ executed in C++ (and in Lean for budgets $40$--$46$)};
\node[lab,below=1pt of thm] {derived in Lean \emph{conditionally} on Step 2};
\end{tikzpicture}
\caption{Logical structure of the proofs of \cref{thm:A,thm:B} and the verification status of each component. \cref{thm:C} has the same structure, with an additional completion routine for two large--large elements (\cref{sec:proof48}); its Step~2 is not stated in Lean.}
\label{fig:status}
\end{figure}

\paragraph{Why large primes are the main difficulty.}
Only finitely many elements of $\Pset$ have both prime factors in a given finite set.
Since $\Pset$ contains $2P$ for every prime $P$, however, no a priori bound on the prime factors of a solution is evident.
The main theoretical step is to show that, once the ``small-prime'' part of a solution is fixed, the large primes are confined to an explicit finite set.
A large prime all of whose neighbours are small must divide an explicit integer $n(A)$ (\cref{lem:star}).
A single edge joining two large primes leads to a factorization identity with finitely many solutions (\cref{prop:component}).
The solution with the prime $3779=29\cdot37+29\cdot41+37\cdot41$ shows that such large primes really occur.

\paragraph{Organization.}
\cref{sec:prelim} gives the size bounds, and \cref{sec:graph} the graph-theoretic form of integrality.
\cref{sec:core,sec:stars} develop the finite reduction.
\cref{sec:proof46} describes the core search, including the loss bound, and proves \cref{thm:A}. \cref{sec:proof47,sec:proof48} prove \cref{thm:B,thm:C}.
\cref{sec:verification,sec:lean,sec:repro} cover verification, formalization and reproducibility.
\cref{sec:discussion} discusses limitations, and \cref{sec:ai} discloses the use of generative AI.

% =====================================================================
\section{Preliminaries and size bounds}\label{sec:prelim}
% =====================================================================

Let $a_1<a_2<a_3<\cdots$ be the elements of $\Pset$ in increasing order, and let
\[
H_k=\sum_{i=1}^{k}\frac1{a_i}\qquad(k\ge 0).
\]
The first $47$ elements of $\Pset$ are
6, 10, 14, 15, 21, 22, 26, 33, 34, 35, 38, 39, 46, 51, 55, 57, 58, 62, 65, 69, 74, 77, 82, 85, 86, 87, 91, 93, 94, 95, 106, 111, 115, 118, 119, 122, 123, 129, 133, 134, 141, 142, 143, 145, 146, 155, 158.

\begin{lemma}[Exchange lemma]\label{lem:exchange}
Let $\mathcal{Q}=\{c_1<c_2<\cdots\}\subset\Z_{>0}$ be infinite. If $T\subset\mathcal{Q}$ is finite with $\card{T}\le k$, then $\Rec{T}\le\sum_{i=1}^{k}1/c_i$. In particular $\Rec{T}\le H_k$ for every $T\subset\Pset$ with $\card{T}\le k$.
\end{lemma}

\begin{proof}
Write $T=\{t_1<\dots<t_m\}$ with $m\le k$. The elements $t_1,\dots,t_i$ are $i$ distinct elements of $\mathcal{Q}$, so $t_i\ge c_i$. Hence $\Rec{T}=\sum_{i\le m}1/t_i\le\sum_{i\le m}1/c_i\le\sum_{i\le k}1/c_i$.
\end{proof}

\begin{proposition}[Exact values]\label{prop:Hk}
The values of $H_k$ in \cref{tab:Hk} are exact reduced fractions. In particular
\[
H_{35}<1-\tfrac{1757}{79000},\qquad H_{37}<1<H_{38},\qquad H_{46}<H_{47}<H_{48}<2 .
\]
\end{proposition}

\begin{table}[ht]
\centering
\caption{Exact values of $H_k$. Each inequality in the last column can be read off by comparing numerator and denominator. The decimal expansions (truncated) are for orientation only.}
\label{tab:Hk}
\small
\setlength{\tabcolsep}{4pt}
\begin{tabular}{@{}lcrrll@{}}
\toprule
 & $a_k$ & numerator & denominator & decimal & \\
\midrule
\tableinput{tables/hk_exact.tex}
\bottomrule
\end{tabular}
\end{table}

\begin{proof}
This is a finite exact computation with rational numbers. It was carried out independently in three ways: with the big-integer routines of the C++ programs (\cref{sec:repro}), in Python with \code{fractions.Fraction} (script \file{paper/verification/verify_paper_facts.py}), and in Lean~4, where the inequalities are proved by \code{norm\_num} (\lean{H35_val}, \lean{H37_lt_one}, \lean{H38_gt_one}, \lean{H46_lt_two}, \lean{H47_lt_two}; \cref{sec:lean}).
Lean also proves that $a_1,\dots,a_{47}$ are the $47$ smallest elements of $\Pset$.
Finally $a_{48}=159$ and $H_{48}=H_{47}+1/159<1.0704$ (checked with exact rationals by \file{paper/verification/verify48_facts.py}; \file{semiprime48.cpp} checks $H_{48}<2$ exactly).
\end{proof}

\begin{corollary}\label{cor:sum-is-one}
Let $T\subset\Pset$ be nonempty with $\card{T}\le 48$ and $\Rec{T}\in\Z$. Then $\Rec{T}=1$ and $\card{T}\ge 38$.
\end{corollary}

\begin{proof}
By \cref{lem:exchange,prop:Hk}, $0<\Rec{T}\le H_{48}<2$, so $\Rec{T}=1$. If $\card{T}\le 37$ then $\Rec{T}\le H_{37}<1$, a contradiction.
\end{proof}

In Lean, the case $\card{T}\le47$ is \lean{recipSum_eq_one_of_integral}. For $\card{T}\le48$ (\lean{recipSum_eq_one_of_integral_48}) Lean uses the cruder bound $\Rec{T}\le H_{47}+1/6<3/2+1/6$, obtained by removing one element.
Consequently, \cref{thm:A} is equivalent to the statement that no set of at most $46$ elements of $\Pset$ has reciprocal sum $1$. Likewise, the $47$-element sets with integral sum in \cref{thm:B}, and the $48$-element sets with integral sum in \cref{thm:C}, are exactly those with sum $1$.

% =====================================================================
\section{The graph of a set of semiprimes and the local criterion}\label{sec:graph}
% =====================================================================

\begin{definition}
For a finite $T\subset\Pset$ let $\Gamma(T)$ be the simple graph whose vertices are the primes dividing some element of $T$ and whose edges are the pairs $\{p,q\}$ with $pq\in T$. For a prime $p$ let $\Nb(p)=\Nb_T(p)=\{q: pq\in T\}$ be its set of neighbours ($\Nb(p)=\varnothing$ if $p$ is not a vertex).
\end{definition}

The map $T\mapsto\Gamma(T)$ is a bijection between finite subsets of $\Pset$ and finite simple graphs on the set of primes without isolated vertices. The number of edges is $\card{T}$. For a prime $p$ and an integer $m$ with $p\nmid m$ we write $m^{-1}$ for the inverse of $m$ modulo $p$.

\begin{lemma}[Local criterion]\label{lem:local}
For every finite $T\subset\Pset$,
\begin{equation}\label{eq:local}
\Rec{T}\in\Z\iff \sum_{q\in\Nb(p)}q^{-1}\equiv 0\pmod p\quad\text{for every prime }p .
\end{equation}
\end{lemma}

\begin{proof}
Let $\ord_p$ denote the $p$-adic valuation on $\Q$. A rational number is an integer if and only if its $p$-adic valuation is nonnegative for every prime $p$.
Fix $p$ and split
\[
\Rec{T}=\frac1p A_p+B_p,\qquad A_p=\sum_{q\in\Nb(p)}\frac1q,\qquad B_p=\sum_{n\in T,\ p\nmid n}\frac1n .
\]
Every element of $T$ is squarefree, so every denominator in $B_p$ and in $A_p$ is prime to $p$, and $\ord_p(B_p)\ge 0$.
Write $A_p=u/v$ with $v=\prod_{q\in\Nb(p)}q$ and $u=\sum_{q\in\Nb(p)}v/q$. Then $p\nmid v$ and $u\equiv v\sum_{q\in\Nb(p)}q^{-1}\pmod p$.
If the congruence in \eqref{eq:local} holds at $p$, then $p\mid u$, so $\ord_p(A_p/p)\ge 0$ and hence $\ord_p(\Rec{T})\ge 0$.
Conversely, if $\ord_p(\Rec{T})\ge0$, then $\ord_p(A_p/p)=\ord_p(\Rec{T}-B_p)\ge 0$, so $p\mid u$ and the congruence holds.
For a prime $p$ that divides no element of $T$ both sides are trivially true.
\end{proof}

\begin{corollary}\label{cor:degree}
If $\Rec{T}\in\Z$, then no vertex of $\Gamma(T)$ has degree $1$.
\end{corollary}

\begin{proof}
If $\Nb(p)=\{q\}$ then $q^{-1}\not\equiv0\pmod p$.
\end{proof}

For every element of the $23$ solutions of \cref{thm:B}, the congruences \eqref{eq:local} were checked directly at every vertex, as an independent test of the implementation.
In Lean the lemma is \lean{integral_iff_localCond}. Its proof uses the common denominator $\prod_{p\in\Gamma(T)}p$ instead of valuations.

% =====================================================================
\section{Core decomposition and the finite reduction}\label{sec:core}
% =====================================================================

Throughout this section $\budget\in\{46,47,48\}$ is a \emph{budget}, and $T\subset\Pset$ is a finite set with
\begin{equation}\label{eq:setting}
\Rec{T}=1,\qquad \card{T}\le\budget .
\end{equation}
By \cref{lem:local}, the congruence \eqref{eq:local} holds at every prime.

\begin{definition}[Small and large primes]
Let $\Sset=\{2,3,5,\dots,73\}$ be the set of the $21$ primes at most $73$. Call the primes $P>73$ \emph{large}; the smallest large primes are $b_0=79<b_1=83<b_2=89<\cdots$.
Let $\Pset_{\Sset}=\{pq\in\Pset: p,q\in\Sset\}$, a set of $\binom{21}{2}=210$ elements.
\end{definition}

The primes dividing $a_1,\dots,a_{46}$ are exactly the elements of $\Sset$ (since $146=2\cdot73\le a_{46}=155<a_{47}=158=2\cdot79$). The reduction below is valid for any finite cutoff; the choice of $\Sset$ only affects the efficiency of the computation.

\begin{definition}[Core, large part, residues]\label{def:core}
For $T$ as in \eqref{eq:setting} define:
\begin{itemize}
\item the \emph{core} $C=T\cap\Pset_{\Sset}$ and the \emph{large part} $G=T\setminus C$;
\item for $q\in\Sset$ the \emph{core residue} $\rho(q)=\sum_{r:\,qr\in C}r^{-1}\in\Z/q\Z$;
\item the \emph{unsatisfied set} $U=U(C)=\{q\in\Sset:\rho(q)\ne0\}$;
\item for $q\in\Sset$ the number $d_q$ of large neighbours of $q$;
\item the number $b$ of \emph{large--large} elements of $G$, i.e.\ elements with both prime factors large.
\end{itemize}
\end{definition}

Every element of $G$ is either $qP$ with $q\in\Sset$ and $P$ large, or $PP'$ with $P,P'$ both large.
In particular every element of $G$ has at most one prime factor in $\Sset$.
The quantities $\rho$ and $U$ depend only on $C$.

\begin{lemma}[Unsatisfied primes]\label{lem:U}
Every $q\in U$ has a large neighbour. Hence $\card{G}\ge\card{U}$ and
\[
\card{C}+\card{U}\le\card{T}\le\budget .
\]
\end{lemma}

\begin{proof}
At $q\in\Sset$ the congruence \eqref{eq:local} reads $\rho(q)+\sum_{P}P^{-1}\equiv0\pmod q$, where $P$ runs over the large neighbours of $q$. If $\rho(q)\ne0$, this sum is nonempty. Choosing one element $e(q)=qP\in G$ for each $q\in U$ gives an injection $U\to G$, because an element of $G$ has at most one prime factor in $\Sset$.
\end{proof}

For $x\ge0$ let $V(x)$ be the sum of the $x$ largest values of $1/n$ over the elements $n\in\Pset$ that have a large prime factor. Thus
$V(x)=\tfrac1{158}+\tfrac1{166}+\tfrac1{178}+\tfrac1{194}+\tfrac1{202}+\tfrac1{206}+\cdots$ ($x$ terms).
Since every such $n$ is at least $2\cdot79=158$, we have $V(x)\le x/158$.

\begin{lemma}[Value bound, general form]\label{lem:value}
For $T$ as in \eqref{eq:setting} we have $\Rec{C}\le1$. Moreover, for \emph{every} subset $U'\subseteq U$,
\begin{equation}\label{eq:value}
1\ \le\ \Rec{C}+\sum_{q\in U'}\frac1{79\,q}+V\bigl(\budget-\card{C}-\card{U'}\bigr).
\end{equation}
\end{lemma}

\begin{proof}
First, $\Rec{C}=1-\Rec{G}\le1$.
Take the elements $e(q)$, $q\in U'$, from the proof of \cref{lem:U}. They are distinct, and $e(q)=qP\ge79q$.
The remaining $\card{G}-\card{U'}\le\budget-\card{C}-\card{U'}$ elements of $G$ are distinct elements of $\Pset$ with a large prime factor.
By \cref{lem:exchange} their reciprocal sum is at most $V(\budget-\card{C}-\card{U'})$.
Now use $1=\Rec{C}+\Rec{G}$.
\end{proof}

\begin{definition}[Admissible cores]\label{def:admissible}
A set $C\subseteq\Pset_{\Sset}$ is \emph{$\budget$-admissible} if $\card{C}+\card{U(C)}\le\budget$ and $\Rec{C}\le1$, and if \eqref{eq:value} holds for every upper segment $U'=U(C)\cap(p,\infty)$, $p\ge0$. (The upper segments include $U(C)$ itself, for $p=0$, and $\varnothing$, for $p\ge73$.)
\end{definition}

By \cref{lem:U,lem:value}, the core of every $T$ satisfying \eqref{eq:setting} is $\budget$-admissible.

\begin{remark}[Why upper segments]\label{rem:segments}
The search of \cref{sec:proof46} processes the primes of $\Sset$ in decreasing order, and at an intermediate stage it uses \eqref{eq:value} with $U'$ equal to the set of primes \emph{already known} to lie in $U$. That set is an upper segment of $U$.
A prime that joins $U$ later, for instance $q=2$ at the last level, adds $1/(79q)$ to the right-hand side of \eqref{eq:value} but removes one term of $V$, possibly a much smaller one. So the right-hand side for $U'=U$ can be larger than for a proper upper segment.
Thus the set of cores produced by the search is \emph{not} the set of cores satisfying \eqref{eq:value} for $U'=U$ alone, as stated in an earlier description of the method. It contains the set of admissible cores in the sense of \cref{def:admissible}, and possibly some further cores, because the bounds used at intermediate nodes are not tight.
Pruning with upper segments is sound because \eqref{eq:value} holds for every $U'\subseteq U$.
On reduced instances the difference is substantial (\cref{sec:bruteforce}).
\end{remark}

\begin{definition}[Shapes]\label{def:shape}
Let $C$ be a core with residues $\rho$ and unsatisfied set $U$. For $q\in U$ let $\Pmin(q)$ be the least large prime $P$ with $P^{-1}\equiv-\rho(q)\pmod q$. Put
\[
w_q(0)=0,\qquad w_q(1)=\frac{1}{q\,\Pmin(q)},\qquad w_q(d)=\sum_{k=0}^{d-1}\frac1{q\,b_k}\quad(d\ge2).
\]
A pair $(d,b)$, with $d=(d_q)_{q\in\Sset}$ and $b$ nonnegative integers, is a \emph{$\budget$-feasible shape} for $C$ if
\begin{enumerate}[label=\textup{(S\arabic*)}]
\item\label{S1} $d_q\ge1$ for $q\in U$ and $d_q\ne1$ for $q\in\Sset\setminus U$;
\item\label{S2} $E:=b+\sum_{q\in U}(d_q-1)+\sum_{q\in\Sset\setminus U}d_q\le\budget-\card{C}-\card{U}$;
\item\label{S3} $\rho(2)+d_2\equiv0\pmod2$, where $\rho(2)\in\{0,1\}$;
\item\label{S4} $1-\Rec{C}\le\sum_{q\in\Sset}w_q(d_q)+\dfrac{b}{79\cdot83}$.
\end{enumerate}
\end{definition}

Dirichlet's theorem guarantees that $\Pmin(q)$ exists, but we do not need it. If no large prime lies in the required class, then $d_q=1$ is impossible by \cref{lem:shape} below. The programs compute $\Pmin(q)$ by direct search, and this search terminated for every class that occurred. The Lean formalization defines $\Pmin(q)=0$ if no such prime exists.

\begin{lemma}[Shape lemma]\label{lem:shape}
For $T$ as in \eqref{eq:setting}, the pair $(d,b)$ of \cref{def:core} is a $\budget$-feasible shape for the core $C$ of $T$.
\end{lemma}

\begin{proof}
\ref{S1}: For $q\in U$ use \cref{lem:U}. If $q\notin U$ has exactly one large neighbour $P$, then \eqref{eq:local} at $q$ gives $P^{-1}\equiv0\pmod q$, which is impossible.

\ref{S2}: Each element of $G$ has exactly one small prime factor or none, so $\card{G}=\sum_{q\in\Sset}d_q+b$. Hence $E=\card{G}-\card{U}\le\budget-\card{C}-\card{U}$.

\ref{S3}: Every large prime $P$ is odd, so $P^{-1}\equiv1\pmod2$. The congruence \eqref{eq:local} at $2$ becomes $\rho(2)+d_2\equiv0\pmod 2$.

\ref{S4}: Split
\[
1-\Rec{C}=\Rec{G}=\sum_{q\in\Sset}\ \sum_{P\text{ large},\,qP\in T}\frac1{qP}+\sum_{PP'\in G\text{ large--large}}\frac1{PP'} .
\]
If $d_q\ge2$, the $d_q$ large neighbours of $q$ are distinct primes, so their contribution is at most $w_q(d_q)$ by \cref{lem:exchange}.
If $d_q=1$, then $q\in U$ by \ref{S1}. By \eqref{eq:local} the unique large neighbour $P$ satisfies $P^{-1}\equiv-\rho(q)$, so $P\ge\Pmin(q)$ and $1/(qP)\le w_q(1)$.
A large--large element is $PP'$ with $79\le P<P'$, so $P'\ge83$ and $1/(PP')\le1/(79\cdot83)$.
\end{proof}

\paragraph{The finite reduction.}
\cref{lem:U,lem:value,lem:shape} reduce the problem as follows.
\begin{enumerate}[label=(\arabic*)]
\item The core is one of finitely many admissible subsets of the $210$ elements of $\Pset_\Sset$.
\item The shape is one of finitely many, because all $d_q$ and $b$ are bounded by \ref{S2}.
\item The large part is then subject to the constraints of \cref{sec:stars}: a large prime whose neighbours are all small divides an explicit positive integer, and a single large--large edge is governed by a factorization identity. Both confine the large primes to explicit finite sets.
\end{enumerate}
The only case not reduced by theory is $b\ge2$. The computation shows that no feasible shape with $b\ge2$ exists for $\budget=47$, and none with $b\ge1$ for $\budget=46$ (\cref{sec:proof46,sec:proof47}). For $\budget=48$, feasible shapes with $b=2$ do occur, and they are handled by the additional argument of \cref{sec:twobb}; no feasible shape with $b\ge3$ exists. The programs were written to report every shape they cannot handle as ``unresolved'' rather than discard it.

% =====================================================================
\section{Large primes: stars and the large--large component}\label{sec:stars}
% =====================================================================

\subsection{Stars}

For a finite set $A$ of primes put
\[
\nA(A)=\sum_{q\in A}\ \prod_{r\in A\setminus\{q\}} r ,
\]
so that $\sum_{q\in A}1/q=\nA(A)/\prod_{r\in A}r$. For instance $\nA(\{a,b\})=a+b$ and $\nA(\{a,b,c\})=ab+ac+bc$.

\begin{lemma}[Star lemma]\label{lem:star}
Let $P$ be a prime with $P\notin A$. Then $\sum_{q\in A}q^{-1}\equiv0\pmod P$ if and only if $P\mid\nA(A)$.
Consequently, if $\Rec{T}\in\Z$ and a large prime $P$ has all its neighbours in $\Sset$, then $P\mid\nA(\Nb(P))$ and in particular $P\le\nA(\Nb(P))$. Moreover $\card{\Nb(P)}\ge2$.
\end{lemma}

\begin{proof}
Multiply the congruence by $\prod_{r\in A}r$, which is invertible modulo $P$. The second statement follows from \eqref{eq:local} at $P$ and $\nA(A)>0$; the third follows from \cref{cor:degree}.
\end{proof}

We call such a pair $(P,A)$ with $A=\Nb(P)\subseteq\Sset$ a \emph{star} with \emph{centre} $P$ and \emph{leaves} $A$.
For a fixed set $A$ of leaves only the finitely many prime divisors of $\nA(A)$ can be centres.
This is the mechanism that bounds the large primes.

\begin{lemma}[Star size]\label{lem:starsize}
Suppose $T$ satisfies \eqref{eq:setting} and $P$ is a large prime all of whose neighbours lie in $\Sset$. Then $\card{\Nb(P)}\le\card{T}-36\le\budget-36$. Thus a star has at most $10$ leaves if $\budget=46$, at most $11$ if $\budget=47$ and at most $12$ if $\budget=48$.
\end{lemma}

\begin{proof}
Let $d=\card{\Nb(P)}$. The $d$ elements $qP$ contribute at most $\frac1{79}\sum_{q\in\Sset}\frac1q\le\frac{1757}{79000}$, since $\sum_{q\in\Sset}1/q=1.7565\ldots\le1.757$ (exact computation).
If $d\ge\card{T}-35$, the other at most $35$ elements contribute at most $H_{35}<1-\frac{1757}{79000}$ by \cref{lem:exchange,prop:Hk}. Then $\Rec{T}<1$, a contradiction.
\end{proof}

The programs use the slightly sharper bound $\frac1{79}\sum_{i<d}1/p_i+H_{\budget-d}\ge1$, with $p_0<p_1<\dots$ the primes of $\Sset$. Recomputed with exact rationals, it gives the same maxima, $10$, $11$ and $12$.

\begin{definition}[Candidate stars]\label{def:candidate}
Let $C$ be a core and $(d,b)$ a feasible shape, and put $W=\{q\in\Sset: d_q\ge1\}$. A \emph{candidate star} for $(C,d)$ with at most $m$ leaves is a pair $(P,A)$ with $A\subseteq W$, $2\le\card{A}\le m$, $P>73$ prime, $P\mid\nA(A)$, and $P^{-1}\equiv-\rho(q)\pmod q$ for every $q\in A$ with $d_q=1$.
\end{definition}

Since $P\le\nA(A)\le\card{A}\prod_{r\in A}r\le 21\prod_{r\in\Sset}r$, there are only finitely many candidate stars (\lean{starCand_finite}).
If $b=0$, every large prime of $T$ is the centre of a star, and the stars of $T$ have the following properties:
\begin{itemize}
\item their centres are distinct;
\item $q\in\Sset$ is a leaf of exactly $d_q$ of them;
\item each of them is a candidate star with at most $\budget-36$ leaves (\cref{lem:star,lem:starsize}). The residue condition holds because a leaf $q$ with $d_q=1$ has the centre as its only large neighbour;
\item their total value $\sum_{(P,A)}\sum_{q\in A}1/(qP)$ equals $1-\Rec{C}$.
\end{itemize}
Conversely, these conditions determine $G$. So for $b=0$ the completions of a core are exactly the \emph{exact covers} of the multiplicities $(d_q)$ by candidate stars with distinct centres and total value $1-\Rec{C}$.

\paragraph{The six solutions with a large prime.}
\cref{tab:bigprimes} lists the six solutions of \cref{thm:B} that contain a large prime.
In each of them $G$ consists of a single star with three leaves $a,b,c$. The core has $44$ elements, $U=\{a,b,c\}$, $d_a=d_b=d_c=1$ and all other $d_q=0$, so the excess is $E=0$ and $\card{C}+\card{U}=47$. The large prime divides $\nA(\{a,b,c\})=ab+ac+bc$.
For $P=311$, $503$ and $3779$ the quantity $ab+ac+bc$ is itself prime and equal to $P$; for $P=269$, $421$ and $433$ it is $3P$, $11P$ and $7P$, respectively.

\begin{table}[ht]
\centering
\caption{The six $47$-term representations of $1$ that contain a prime $P>73$. The last column refers to the numbering of \cref{app:solutions}.}
\label{tab:bigprimes}
\small
\begin{tabular}{@{}rllll@{}}
\toprule
$P$ & leaves $A=\{a,b,c\}$ & $\nA(A)=ab+ac+bc$ & large elements $aP,bP,cP$ & no. \\
\midrule
\tableinput{tables/bigprime_table.tex}
\bottomrule
\end{tabular}
\end{table}

\begin{example}[The prime $3779$]\label{ex:3779}
Solution~5 of \cref{app:solutions} contains $109591=29\cdot3779$, $139823=37\cdot3779$ and $154939=41\cdot3779$, and $3779=29\cdot37+29\cdot41+37\cdot41$ is prime.
Any search that bounds the prime factors by a fixed cutoff below $3779$ misses this representation of $1$. Such a cutoff includes the natural choices $101$ and $1000$.
The present argument bounds large primes only through divisibility: a star centre divides $\nA(A)$ for its set of leaves $A$, which for three leaves can be as large as $67\cdot71+67\cdot73+71\cdot73=14\,831$.
\end{example}

\subsection{One edge between two large primes}\label{sec:component}

Suppose $b=1$. Then $G$ contains exactly one element $P_1P_2$ with $P_1\ne P_2$ both large.
Every other large prime has only small neighbours, since a second large neighbour would create a second large--large element, and is therefore the centre of a star.
For $i=1,2$ we have $\Nb(P_i)=A_i\cup\{P_{3-i}\}$ with $A_i\subseteq\Sset$, and $A_i\ne\varnothing$ by \cref{cor:degree}.
Put
\[
b_i=\prod_{q\in A_i}q,\qquad a_i=\nA(A_i),\qquad Z=\sum_{q\in A_1}\frac1{qP_1}+\sum_{q\in A_2}\frac1{qP_2}+\frac1{P_1P_2},
\]
the value of the component. Finally put
\[
W=a_1b_2P_2+a_2b_1P_1+b_1b_2,\qquad t=Z\,b_1b_2,\qquad N=b_1b_2\,(t+a_1a_2).
\]

\begin{proposition}[The component equation]\label{prop:component}
Let $P_1\ne P_2$ be primes not dividing $b_1b_2$, and let $A_1,A_2$ be nonempty.
\begin{enumerate}[label=\textup{(\alph*)}]
\item $Z\,b_1b_2\,P_1P_2=W$; that is, $t\,P_1P_2=W$ in $\Q$.
\item The local conditions \eqref{eq:local} at $P_1$ and at $P_2$ hold if and only if $P_1P_2\mid W$, i.e.\ if and only if $t\in\Z$.
\item For every rational $t\ne0$,
\[
(tP_1-a_1b_2)(tP_2-a_2b_1)=b_1b_2\,(t+a_1a_2)\iff tP_1P_2=W .
\]
In particular the identity on the left holds for $t=Zb_1b_2$, whether or not $t$ is an integer.
\item Suppose $t$ is a positive integer with $tP_1P_2=W$, and put $X=tP_1-a_1b_2$, $Y=tP_2-a_2b_1$. Then $X$ and $Y$ are positive integers with $XY=N$, and
\[
P_1=\frac{X+a_1b_2}{t},\qquad P_2=\frac{N/X+a_2b_1}{t}.
\]
Conversely, if $t>0$ and $X$ is a positive divisor of $N$ such that both quotients are integers, then they satisfy $tP_1P_2=W$.
\end{enumerate}
\end{proposition}

\begin{proof}
(a) Use $\sum_{q\in A_i}1/q=a_i/b_i$ and clear denominators.

(b) Since $\Nb(P_1)=A_1\cup\{P_2\}$, \cref{lem:star} shows that \eqref{eq:local} at $P_1$ is equivalent to $P_1\mid\nA(A_1\cup\{P_2\})=a_1P_2+b_1$.
Now $W=b_2(a_1P_2+b_1)+a_2b_1P_1\equiv b_2(a_1P_2+b_1)\pmod{P_1}$ and $P_1\nmid b_2$, so this is equivalent to $P_1\mid W$. The same argument applies at $P_2$.
Since $P_1\ne P_2$, both conditions together are equivalent to $P_1P_2\mid W$, which by (a) says $t\in\Z$.

(c) Expanding gives
$(tP_1-a_1b_2)(tP_2-a_2b_1)-b_1b_2(t+a_1a_2)=t\,(tP_1P_2-W)$.

(d) From $tP_1P_2=W$ we get $P_2X=W-a_1b_2P_2=a_2b_1P_1+b_1b_2>0$, so $X>0$, and similarly $Y>0$. Then (c) gives $XY=N$. The converse is (c) read backwards.
\end{proof}

\begin{remark}[Correction]\label{rem:component}
An earlier description of the method stated that the factorization identity in (c) is itself equivalent to the local conditions at $P_1$ and $P_2$. By (a) and (c), however, the identity holds automatically for $t=Zb_1b_2$. What is equivalent to the local conditions is the \emph{integrality} of $t$, by (b). The positivity of $X$ and $Y$ and the nonemptiness of $A_1$ and $A_2$, both used by the program, are proved above.
\end{remark}

\paragraph{Finite enumeration for $b=1$.}
Fix a core $C$, a feasible shape with $b=1$, nonempty attachment sets $A_1,A_2\subseteq W$, and a family $F$ of stars covering the remaining multiplicities $d_q-[q\in A_1]-[q\in A_2]$.
Then the value of the component is $Z=1-\Rec{C}-\Rec{F}$, and $t=Zb_1b_2$ is determined; it must be a positive integer.
By \cref{prop:component}(d) the pair $(P_1,P_2)$ is then determined by a positive divisor $X$ of the explicit integer $N$.
All choices of $A_1$, $A_2$ and $F$ range over finite sets, so the $b=1$ completions of a core form a finite, explicitly enumerable set (\lean{comp_pairs_finite}).
Interchanging $(A_1,P_1)$ and $(A_2,P_2)$ replaces $X$ by $N/X$, so unordered pairs $\{A_1,A_2\}$ suffice.

% =====================================================================
\section{The core search and the proof of the lower bound}\label{sec:proof46}
% =====================================================================

This section describes the computation and proves \cref{thm:A} from its outcome for $\budget=46$.
The computation consists of two steps: enumeration of all admissible cores (Step~1) and, for each of them, enumeration of all feasible shapes and completions (Step~2).
Both steps are carried out by the program \file{semiprime48.cpp} for any budget $\budget\le48$.
Its Step~1 is the depth-first search of the earlier programs \file{semiprime_reciprocals.cpp} (budget $46$) and \file{semiprime47.cpp} (budget $47$), with one additional pruning test, the loss bound (P4) of \cref{sec:lossbound}. That test reduces the number of search nodes by a factor of about $5000$ for budget $46$ and $3000$ for budget $47$, and it does not change the output (\cref{lem:sameoutput}).

\subsection{Step 1: enumeration of the admissible cores}\label{sec:step1}

Let $p_0=2<p_1<\dots<p_{20}=73$ be the primes of $\Sset$.
Every core $C\subseteq\Pset_\Sset$ is determined by the sets
\[
L_j=\{p_i: i<j,\ p_ip_j\in C\}\qquad(j=20,19,\dots,1)
\]
of \emph{lower neighbours}. The search is a depth-first search that chooses $L_{20},L_{19},\dots,L_1$ in this order (\cref{alg:step1}).
When $L_j$ is chosen, every core edge at $p_j$ is known, because the edges to larger primes were fixed at earlier levels. Hence $\rho(p_j)$ is known, and it is decided whether $p_j\in U$.
At every stage the set of primes already known to lie in $U$ is therefore an upper segment $U\cap(p,\infty)$ of the final set $U$.
Each core corresponds to exactly one leaf of the search tree, so each core is visited at most once.

For each $j$, the choice of $L_j$ is split into two parts. The primes $p_i$ with $8\le i<j$ are handled by explicit inclusion/exclusion recursion. The $8$ smallest primes are handled by a precomputed table of all $2^8$ subsets, indexed by their contribution to the residue at $p_j$. This lets the search jump directly to the subsets that make $p_j$ satisfied, and to the others (which put $p_j$ into $U$) only when that is still affordable.

\paragraph{Fixed-point bounds.}
Write $\up(n)=\lceil 2^{96}/n\rceil$ and $\lo(n)=\lfloor 2^{96}/n\rfloor$. For a finite set $X$ of positive integers put $\up(X)=\sum_{n\in X}\up(n)$ and $\lo(X)=\sum_{n\in X}\lo(n)$. Then
\[
\lo(X)\le 2^{96}\Rec{X}\le\up(X).
\]
Let $v_1<v_2<\cdots$ be the elements of $\Pset$ with a large prime factor, and put $\mathrm{VB}(x)=\sum_{i\le x}\up(v_i)\ge2^{96}V(x)$.
All these quantities are integers below $2^{128}$ and are computed exactly in unsigned $128$-bit arithmetic.

\paragraph{Pruning tests.}
Consider a node of the search at which the decided core edges form a set $C_0$ and the decided unsatisfied primes form $U_0$. Let $j$ be the current level and $i$ the number of lower primes whose membership in $L_j$ is still open. The search discards the node if one of the following holds:
\begin{enumerate}[label=\textup{(P\arabic*)}]
\item\label{P1} $\card{C_0}+\card{U_0}>\budget$;
\item\label{P2} $\lo(C_0)>2^{96}$;
\item\label{P3} $\up(C_0)+\sum_{q\in U_0}\up(79q)+B_{j,i}(R)<2^{96}$, where $R=\budget-\card{C_0}-\card{U_0}$ and
\[
B_{j,i}(R)=\max_{\substack{x_1+x_2\le R\\ x_1\le i}}\Bigl(M_{j,i}(x_1)+M'_j(x_2)+\mathrm{VB}(R-x_1-x_2)\Bigr).
\]
Here $M_{j,i}(x)$ is the sum of the $x$ largest values $\up(p_jp_k)$ with $k<i$ (possible further edges at $p_j$), and $M'_j(x)$ is the sum of the $x$ largest values $\up(p_kp_l)$ with $k<l<j$ (possible edges among smaller primes).
\end{enumerate}
Test \ref{P3} is also applied to the tentative set $U_0\cup\{p_j\}$, with $R$ replaced by $R-1$, before $p_j$ is put into $U$.
The fourth test, \ref{P4}, is described in \cref{sec:lossbound}.
At a leaf, where the complete core $C$ and $U$ are known, the core is kept only if \ref{P1} and \ref{P2} hold and
\begin{equation}\label{eq:leaftest}
\mathrm{Val}(C):=\up(C)+\sum_{q\in U}\up(79q)+\mathrm{VB}(\budget-\card{C}-\card{U})\ge2^{96}.
\end{equation}

\begin{algorithm}[t]
\caption{Step 1 (core enumeration), schematically}\label{alg:step1}
\begin{algorithmic}[1]
\Procedure{Level}{$j$, $C_0$, $U_0$} \Comment{primes $>p_j$ fully processed}
  \If{$j=0$} \Comment{prime $2$: no lower neighbours}
     \State decide $2\in U$ from $\rho(2)$; apply \ref{P1}; \Call{Leaf}{$C_0$, $U_0$}; \Return
  \EndIf
  \ForAll{$L\subseteq\{p_0,\dots,p_{j-1}\}$ not excluded by \ref{P1}--\ref{P4}} \Comment{explicit + table}
     \State $C_1\gets C_0\cup\{p_jq:q\in L\}$; compute $\rho(p_j)$ and the residue updates at $q\in L$
     \State $U_1\gets U_0\cup\{p_j\}$ if $\rho(p_j)\ne0$ else $U_0$ \Comment{\ref{P3} tested for $U_1$}
     \State \Call{Level}{$j-1$, $C_1$, $U_1$}
  \EndFor
\EndProcedure
\Procedure{Leaf}{$C$, $U$}
  \If{\eqref{eq:leaftest} holds} output $C$ to Step~2 \EndIf
\EndProcedure
\end{algorithmic}
\end{algorithm}

\subsection{The loss bound}\label{sec:lossbound}

\paragraph{Why the search without \ref{P4} is slow.}
Profiling the search with \ref{P1}--\ref{P3} only shows that almost all nodes lie in branches in which some prime has been put into $U$.
For instance, for budget $45$ the search visits $1.07\cdot10^{10}$ nodes, but only $2.9\cdot10^{7}$ when unsatisfied primes are switched off (the latter is not a valid search; it only shows where the nodes are).
Once a prime is in $U$, every set of neighbours is allowed at it, and the value bound \ref{P3} ignores the congruences \eqref{eq:local} that the primes not yet processed will have to satisfy. Satisfying them forces a loss of value that \ref{P3} cannot see, so such branches die only many levels further down.
Test \ref{P4} makes this loss visible, prime by prime.

\paragraph{Weights.}
All quantities are integers in units of $2^{-96}$. Fix $\theta=\lfloor2^{96}/145\rfloor$ (any $\theta\ge0$ would be valid; this value was tuned by experiment). For an edge $e=p_ap_b$ with $a<b$ and for $q\in\Sset$ put
\[
w(e)=\up(p_ap_b)-\theta,\qquad g(q)=\up(79q)-\theta,\qquad V_\theta=\max_{0\le x\le\budget}\bigl(\mathrm{VB}(x)-x\theta\bigr).
\]
If $C$ is a core with $\card{C}+\card{U}\le\budget$ and $x=\budget-\card{C}-\card{U}$, then
\begin{equation}\label{eq:lagrange}
\mathrm{Val}(C)=\budget\theta+\sum_{e\in C}w(e)+\sum_{q\in U}g(q)+\bigl(\mathrm{VB}(x)-x\theta\bigr)\ \le\ \budget\theta+V_\theta+\sum_{e\in C}w(e)+\sum_{q\in U}g(q).
\end{equation}
The weights may be negative; $\theta$ acts as a Lagrange multiplier for the budget constraint.

\paragraph{Shares.}
The weight of an edge $e=p_ap_b$, $a<b$, is split between its endpoints. With $\sigma_a=s_a/12$, where $(s_0,\dots,s_5)=(12,11,11,10,10,9)$ and $s_a=8$ for $a\ge6$, the larger endpoint receives $\mathrm{sh}_b(e)=\lceil\sigma_a w(e)\rceil$ and the smaller one $\mathrm{sh}_a(e)=\lceil(1-\sigma_a)w(e)\rceil$. Then $\mathrm{sh}_a(e)+\mathrm{sh}_b(e)\ge w(e)$. (Any $\sigma_a\in[0,1]$ is valid; these values were tuned by experiment.)

\paragraph{Local bounds.}
Consider a prime $p_k$ whose membership in $U$ is not yet decided, let $X$ be the set of its possible further neighbours, and let $r$ be the residue at $p_k$ of its decided edges. In a final core, either $p_k\in U$, or the set $Y\subseteq X$ of its further neighbours satisfies $r+\sum_{y\in Y}y^{-1}\equiv0\pmod{p_k}$. Hence the shares received by $p_k$ from its further edges, plus $g(p_k)$ if $p_k\in U$, total at most
\[
\mathrm{loc}(k;X,r)=\max\Bigl(g(p_k)+\sum_{x\in X}\max\bigl(\mathrm{sh}_k(p_kx),0\bigr),\ \max_{\substack{Y\subseteq X\\ r+\sum_{y\in Y}y^{-1}\equiv0\ (p_k)}}\ \sum_{y\in Y}\mathrm{sh}_k(p_ky)\Bigr).
\]
The inner maximum is a knapsack over the residues modulo $p_k$. At a node of the search, $X$ depends only on the level $(j,i)$ and on $k$, so every value of $\mathrm{loc}$ that can occur is a precomputed table entry indexed by the level, $k$ and $r$ (about $20\,000$ entries).

\begin{enumerate}[label=\textup{(P\arabic*)},start=4]
\item\label{P4} At a node with decided core edges $C_0$ and decided unsatisfied primes $U_0$, the search discards the node if
\[
\up(C_0)+\sum_{q\in U_0}\up(79q)+\bigl(\budget-\card{C_0}-\card{U_0}\bigr)\theta+V_\theta+\sum_{k}\mathrm{loc}(k;X_k,r_k)<2^{96},
\]
where $k$ runs over the primes not yet decided, $X_k$ is the set of their possible further neighbours and $r_k$ the residue of their decided edges.
\end{enumerate}

The sum over $k$ changes by a bounded amount when an edge is added or rejected, and for the table-driven step a $256$-entry array gives it for all subsets of the $8$ smallest primes at once.
All quantities are integers, and the shares are rounded up.

\begin{proposition}[Soundness of the loss bound]\label{prop:loss}
For every core $C$ with $\card{C}+\card{U(C)}\le\budget$ in the subtree of a node, the left-hand side of \ref{P4} is at least $\mathrm{Val}(C)$. In particular \ref{P4} never discards a node above a core that passes \eqref{eq:leaftest}.
\end{proposition}

\begin{proof}
Split $C=C_0\cup C_1$ and $U=U_0\cup U_1$ into the decided and undecided parts. Then $\sum_{C_0}w+\sum_{U_0}g=\up(C_0)+\sum_{U_0}\up(79q)-(\card{C_0}+\card{U_0})\theta$, so by \eqref{eq:lagrange} it suffices to show $F:=\sum_{e\in C_1}w(e)+\sum_{q\in U_1}g(q)\le\sum_k\mathrm{loc}(k;X_k,r_k)$.
Every edge of $C_1$ joins two undecided primes, and $w(e)\le\mathrm{sh}_a(e)+\mathrm{sh}_b(e)$. Hence $F$ is at most the sum over the undecided primes $p_k$ of the shares that $p_k$ receives from its edges in $C_1$, plus $g(p_k)$ if $p_k\in U_1$. By the choice of $\mathrm{loc}$, the contribution of $p_k$ is at most $\mathrm{loc}(k;X_k,r_k)$.
\end{proof}

\begin{lemma}[Soundness of Step 1]\label{lem:step1}
Every $\budget$-admissible core (\cref{def:admissible}) reaches a leaf of the search and is passed to Step~2.
\end{lemma}

\begin{proof}
Let $C$ be admissible, with unsatisfied set $U$, and consider a node on the path to $C$, with $C_0\subseteq C$ and $U_0=U\cap(p,\infty)$ for some $p$.
\ref{P1} fails because $\card{C_0}+\card{U_0}\le\card{C}+\card{U}\le\budget$.
\ref{P2} fails because $\lo(C_0)\le2^{96}\Rec{C_0}\le2^{96}\Rec{C}\le2^{96}$.
For \ref{P3}, let $x_1$ and $x_2$ be the numbers of edges of $C\setminus C_0$ at $p_j$ and among the smaller primes, respectively. Then $\up(C\setminus C_0)\le M_{j,i}(x_1)+M'_j(x_2)$. Also $R-x_1-x_2=\budget-\card{C}-\card{U_0}$, and $\mathrm{VB}(x)\ge2^{96}V(x)$. Hence
\[
\up(C_0)+\sum_{q\in U_0}\up(79q)+B_{j,i}(R)\ \ge\ 2^{96}\Bigl(\Rec{C}+\sum_{q\in U_0}\frac1{79q}+V(\budget-\card{C}-\card{U_0})\Bigr)\ \ge\ 2^{96},
\]
where the last step is \eqref{eq:value} for the upper segment $U'=U_0$.
The tentative test with $U_0\cup\{p_j\}$ is again a test of this form for the upper segment $U\cap[p_j,\infty)$, and \eqref{eq:leaftest} is the case $U'=U$; in particular $\mathrm{Val}(C)\ge2^{96}$.
Finally \ref{P4} fails by \cref{prop:loss}.
\end{proof}

Because of the rounding, the search may keep a few cores that are not admissible. This is harmless: they are examined further in Step~2.
\cref{lem:step1} is proved in Lean for the search written as a Lean function (\lean{core_mem_search}, \cref{sec:lean}).

\begin{lemma}[Same output]\label{lem:sameoutput}
The search with \ref{P1}--\ref{P4} outputs exactly the same cores as the search with \ref{P1}--\ref{P3} only.
\end{lemma}

\begin{proof}
Adding a test can only remove cores. Conversely, a core output by the search without \ref{P4} passes \eqref{eq:leaftest}, so by \cref{prop:loss} test \ref{P4} does not discard any node above it.
\end{proof}

In agreement with \cref{lem:sameoutput}, the archived core lists of the earlier programs for budgets $46$ and $47$ are reproduced line for line by \file{semiprime48.cpp} (\cref{sec:rerun}).

\paragraph{Parallelization.}
The search tree above a fixed level (the prime $43$ in \file{semiprime48.cpp}, $61$ in the earlier programs) is traversed by every thread in the same deterministic order, and the nodes at that level are numbered consecutively.
Each thread obtains node numbers from a shared atomic counter and processes only the nodes whose numbers it obtained.
Tickets are handed out in increasing order, and a thread requests a new ticket only after passing the node of its previous ticket. Hence the ticket a thread obtains is never smaller than its current position, and every node number below the total count is obtained by exactly one thread.
So every subtree is processed exactly once, regardless of scheduling. Only the reported node counts depend on the number of threads, because the upper part of the tree is traversed by every thread.

\subsection{Step 2: completion of the surviving cores}\label{sec:step2}

For each core $C$ output by Step~1:
\begin{enumerate}[label=(\alph*)]
\item If $U=\varnothing$, then $C$ satisfies \eqref{eq:local} at every prime: at the primes of $\Sset$ by the definition of $U$, and trivially at all other primes. So $\Rec{C}\in\Z$. If $C\ne\varnothing$, then $\Rec{C}=1$ by \cref{cor:sum-is-one}, and the program reports $C$ as a solution. The empty core is not admissible, because $V(\budget)\le\budget/158<1$.
\item If $U\ne\varnothing$, the program enumerates all $(d,b)$ satisfying \ref{S1} and \ref{S2}, discards those that violate \ref{S3} or \ref{S4}, and processes the survivors.
In \ref{S4} the right-hand side is rounded up and $1-\Rec{C}$ is rounded down, so no feasible shape is discarded.
For a surviving shape with $b=0$, the program lists all candidate stars with at most $\budget-36$ leaves (\cref{def:candidate}); it finds the centres by factoring $\nA(A)$ completely.
It then enumerates the exact covers of the multiplicities by candidate stars with distinct centres whose total value is $1-\Rec{C}$.
Surviving shapes with $b=1$ and $b=2$ are treated in \cref{sec:proof47,sec:twobb}; a surviving shape with $b\ge3$ is reported as \emph{unresolved}.
\end{enumerate}
For budget $46$ the earlier program \file{semiprime_reciprocals.cpp} only searched for one cover satisfying \eqref{eq:local} at every $q\in W$ and reported every shape with $b>0$ as unresolved; for this budget the difference does not matter, because no candidate star and no shape with $b>0$ occurs.

\begin{claim}[Budget $46$; Lean: \protect\lean{Search46}, \protect\lean{Search46'}]\label{claim:46}
For every $46$-admissible core $C$:
\begin{enumerate}[label=\textup{(\roman*)}]
\item $U(C)\ne\varnothing$;
\item every $46$-feasible shape $(d,b)$ for $C$ has $b=0$;
\item for every $46$-feasible shape $(d,0)$ there is no candidate star with at most $10$ leaves.
\end{enumerate}
\end{claim}

\begin{table}[t]
\centering
\caption{Recorded results of \protect\file{semiprime48.cpp} with $4$ threads (\protect\file{validation48/run46_output.txt}, \protect\file{validation48/run47_output.txt}, \protect\file{run48_output.txt}). For budgets $46$ and $47$ the earlier programs (\protect\file{proof_run_output.txt}, \protect\file{run47_output.txt}) printed the same values for all quantities they report; they did not report the split by $b$ (for budget $47$ it was obtained at the time with the Python program of \cref{sec:python}, with the same result). Their node counts and times are given in the last rows.}
\label{tab:results}
\small
\begin{tabular}{@{}lrrr@{}}
\toprule
 & budget $46$ & budget $47$ & budget $48$\\
\midrule
maximal number of leaves of a star (\cref{lem:starsize}) & $10$ & $11$ & $12$\\
search nodes & $26\,112\,391$ & $598\,305\,041$ & $13\,344\,629\,391$\\
cores kept by Step 1 & $178$ & $22\,382$ & $1\,491\,334$\\
\quad of which $U=\varnothing$ & $0$ & $17$ & $414$\\
\quad of which $U\ne\varnothing$ and $\card{C}+\card{U}=\budget$ & $111$ & $11\,287$ & $702\,715$\\
shapes satisfying \ref{S1}--\ref{S3} & $3\,777$ & $1\,107\,702$ & $116\,874\,859$\\
shapes satisfying \ref{S1}--\ref{S4} & $61$ & $8\,049$ & $601\,685$\\
\quad cores with at least one such shape & $61$ & $7\,432$ & $509\,055$\\
\quad with $b=0$ & $61$ & $7\,981$ & $592\,627$\\
\quad with $b=1$ & $0$ & $68$ & $8\,985$\\
\quad with $b=2$ & $0$ & $0$ & $73$\\
\quad with $b\ge3$ & $0$ & $0$ & $0$\\
candidate stars (summed over surviving shapes) & $0$ & $125$ & $24\,065$\\
unresolved cores & $0$ & $0$ & $0$\\
probable (not proven) primes used & $0$ & $0$ & $0$\\
solutions with $\card{T}=\budget$ & $0$ & $23$ & $620$\\
wall-clock time (search and completion) & $2.1$\,s & $65$\,s & $29$\,min\,$30$\,s\\
\midrule
search nodes without \ref{P4} (earlier programs) & $139\,830\,520\,211$ & $1\,795\,181\,713\,099$ & --\\
wall-clock time of the earlier programs & 19\,min\,40\,s & 5\,h\,48\,min & --\\
\bottomrule
\end{tabular}
\end{table}

The run of \file{semiprime48.cpp} with budget $46$ produced the output summarized in \cref{tab:results}. The search visited $2.6\cdot10^{7}$ nodes and kept $178$ cores; none of them has $U=\varnothing$. Step~2 examined $3777$ shapes satisfying \ref{S1}--\ref{S3}. Only $61$ of them also satisfy \ref{S4}, all with $b=0$, and for none of these $61$ shapes does a candidate star exist. Nothing was reported as unresolved and no solution was found.
This is precisely \cref{claim:46}; see \cref{sec:claimmatch} for the correspondence in detail.
The recorded run of the earlier program \file{semiprime_reciprocals.cpp}, which visited $1.4\cdot10^{11}$ nodes, kept the same $178$ cores and printed the same Step~2 counts; it was repeated for the first version of this paper with identical results (\cref{sec:rerun}).

\begin{proof}[Proof of \cref{thm:A}, assuming \cref{claim:46}]
Let $T\subset\Pset$ be nonempty with $\card{T}\le46$ and $\Rec{T}\in\Z$.
By \cref{cor:sum-is-one}, $\Rec{T}=1$, so $T$ satisfies \eqref{eq:setting} with $\budget=46$.
Its core $C$ is $46$-admissible (\cref{lem:U,lem:value}), so $U\ne\varnothing$ by \cref{claim:46}(i).
The shape $(d,b)$ of $T$ is $46$-feasible (\cref{lem:shape}), so $b=0$ by \cref{claim:46}(ii).
Choose $q\in U$. By \cref{lem:U}, $q$ has a large neighbour $P$. Since $b=0$, all neighbours of $P$ lie in $\Sset$.
By \cref{lem:star,lem:starsize}, $(P,\Nb(P))$ satisfies $2\le\card{\Nb(P)}\le10$ and $P\mid\nA(\Nb(P))$, and $\Nb(P)\subseteq W$. Moreover, if $d_{q'}=1$ for some $q'\in\Nb(P)$, then $P$ is the only large neighbour of $q'$, so $P^{-1}\equiv-\rho(q')\pmod{q'}$ by \eqref{eq:local}.
Thus $(P,\Nb(P))$ is a candidate star for $(C,d)$, contradicting \cref{claim:46}(iii).
Hence $\card{T}\ge47$. Sharpness follows from \cref{thm:B}, or directly from any of the sets in \cref{app:solutions}, whose validity is verified exactly and unconditionally (\cref{sec:exact}).
\end{proof}

% =====================================================================
\section{Classification of the 47-term representations}\label{sec:proof47}
% =====================================================================

For $\budget=47$ the program \file{semiprime48.cpp} runs Step~1 as in \cref{sec:proof46}. Its Step~2 enumerates \emph{all} completions, as did Step~2 of the earlier program \file{semiprime47.cpp}:
\begin{itemize}
\item cores with $U=\varnothing$ are recorded as solutions;
\item for a surviving shape with $b=0$, all exact covers by candidate stars with at most $11$ leaves and distinct centres are enumerated, and a cover is kept if its value equals $1-\Rec{C}$ \emph{exactly}. The value is computed as an integer multiple of $1/D$, where $D=\prod_{p\in\Sset}p<2^{128}$ is a common denominator of all core and star values (the value of a star $(P,A)$ is $(\nA(A)/P)/\prod_{q\in A}q$);
\item for a surviving shape with $b=1$, the program enumerates the following, as in \cref{sec:component}:
  \begin{itemize}
  \item all unordered pairs $\{A_1,A_2\}$ of nonempty subsets of $W$ with $[q\in A_1]+[q\in A_2]\le d_q$;
  \item all covers of the remaining multiplicities by candidate stars;
  \item the integer $t$, if it exists;
  \item all positive divisors $X$ of $N$, obtained from the complete factorization of $t+a_1a_2$ and of $b_1b_2$.
  \end{itemize}
  It keeps the pairs for which $P_1\ne P_2$ are primes larger than $73$ and distinct from the star centres;
\item a surviving shape with $b=2$ is treated as in \cref{sec:twobb} (in \file{semiprime47.cpp}: reported as unresolved), and a surviving shape with $b\ge3$ is reported as unresolved.
\end{itemize}
Every recorded set is then verified exactly: it has $47$ elements, they are distinct squarefree semiprimes, and their reciprocal sum is $1$ (\cref{sec:exact}).

\begin{claim}[Budget $47$; Lean: \protect\lean{Search47}, \protect\lean{Search47'}]\label{claim:47}
Let $\mathcal{T}_{23}$ be the family of the $23$ sets listed in \cref{app:solutions}.
\begin{enumerate}[label=\textup{(\roman*)}]
\item Every $47$-admissible core $C\ne\varnothing$ with $U(C)=\varnothing$ belongs to $\mathcal{T}_{23}$.
\item For every $47$-admissible core with $U\ne\varnothing$, every $47$-feasible shape has $b\le1$.
\item \textup{($b=0$)} Let $C$ be $47$-admissible with $U\ne\varnothing$, let $(d,0)$ be a $47$-feasible shape, and let $\{(P,A_P)\}_{P\in\mathcal{C}}$ be a family of candidate stars with at most $11$ leaves and distinct centres such that each $q\in\Sset$ is a leaf of exactly $d_q$ of them and $\Rec{C}+\sum_{P}\sum_{q\in A_P}1/(qP)=1$. Then $C\cup\{qP: P\in\mathcal{C},\,q\in A_P\}\in\mathcal{T}_{23}$.
\item \textup{($b=1$)} Let $C$ and $(d,1)$ be as in (iii), except that $b=1$. Suppose we are given:
  \begin{itemize}
  \item nonempty $A_1,A_2\subseteq\Sset$;
  \item a family of candidate stars as in (iii) covering the multiplicities $d_q-[q\in A_1]-[q\in A_2]$, with total value $\Sigma_F<1-\Rec{C}$;
  \item the positive integer $t=(1-\Rec{C}-\Sigma_F)\,b_1b_2$;
  \item a positive divisor $X$ of $N$ such that $P_1=(X+a_1b_2)/t$ and $P_2=(N/X+a_2b_1)/t$ are distinct primes larger than $73$ and are not centres of the stars.
  \end{itemize}
  Then the resulting set, consisting of $C$, the star elements, $\{qP_1\}_{q\in A_1}$, $\{qP_2\}_{q\in A_2}$ and $P_1P_2$, belongs to $\mathcal{T}_{23}$.
\end{enumerate}
\end{claim}

The run of \file{semiprime48.cpp} and the recorded run of \file{semiprime47.cpp} (\cref{tab:results}) kept the same $22\,382$ cores. Seventeen of them have $U=\varnothing$, all with $\card{C}=47$, and these are the $17$ solutions with all primes at most $73$.
Of the remaining $22\,365$ cores, $11\,287$ have $\card{C}+\card{U}=47$.
The printed count ``$\card{C}+\card{U}=47$: $11\,287$'' in \file{run47_output.txt} excludes the $17$ cores with $U=\varnothing$. Including them, the dump of all kept cores contains $11\,304$ cores with $\card{C}+\card{U}=47$.
Step~2 examined $1\,107\,702$ shapes. Of these, $8\,049$ shapes in $7\,432$ cores satisfy \ref{S1}--\ref{S4}: $7\,981$ with $b=0$, $68$ with $b=1$, and none with $b\ge2$.
For the $68$ shapes with $b=1$ there are $7\,896$ admissible attachment pairs; $t$ is an integer in $957$ cases, and none of them yields a solution.
The surviving shapes with $b=0$ have $125$ candidate stars in total. They produce exactly the $6$ solutions of \cref{tab:bigprimes}, each consisting of a core with $44$ elements and a single star with three leaves.
No shape was unresolved. In the run of \file{semiprime47.cpp} every number declared prime is below $2^{64}$, where the primality test is deterministic, and ``composite'' verdicts of the Miller--Rabin test are always correct; \file{semiprime48.cpp} reports $0$ probable primes. So all primality decisions are rigorous (\cref{sec:repro}).
Both programs found exactly the same $23$ solutions.

\begin{proof}[Proof of \cref{thm:B}, assuming \cref{claim:47}]
By \cref{sec:exact}, each of the $23$ listed sets has $47$ elements and reciprocal sum $1$, and the sets are pairwise distinct.
Conversely, let $\card{T}=47$ and $\Rec{T}=1$. Then \eqref{eq:setting} holds with $\budget=47$, and the core $C$ of $T$ is $47$-admissible.
\begin{itemize}
\item If $U=\varnothing$, then $G=\varnothing$. Indeed, $C$ satisfies \eqref{eq:local} everywhere, so $\Rec{C}\in\Z$. Since $V(47)<1$, the core $C$ is nonempty, so $\Rec{C}=1$ and $\Rec{G}=0$. Hence $T=C\in\mathcal{T}_{23}$ by \cref{claim:47}(i).
\item If $U\ne\varnothing$, the shape of $T$ is feasible, so $b\le1$ by \cref{claim:47}(ii).
  \begin{itemize}
  \item If $b=0$, the stars of $T$ form a family as in \cref{claim:47}(iii) (\cref{sec:stars}), so $T\in\mathcal{T}_{23}$.
  \item If $b=1$, \cref{prop:component} shows that the data $(A_1,A_2,\text{stars},t,X)$ of $T$ are as in \cref{claim:47}(iv), with $X=tP_1-a_1b_2>0$. Hence $T\in\mathcal{T}_{23}$.
  \end{itemize}
\end{itemize}
The statement about integral sums follows from \cref{cor:sum-is-one}. Parts (i) and (ii) of \cref{thm:B} are read off from the list; the numbers in \cref{tab:bigprimes} were recomputed exactly.
\end{proof}

\paragraph{Structure of the $23$ solutions.}
All $23$ solutions contain the $23$ numbers
6, 10, 14, 15, 21, 22, 26, 33, 34, 35, 38, 39, 46, 51, 55, 57, 58, 62, 65, 93, 95, 133, 155.
Each solution involves between $14$ and $16$ distinct primes, and the union of the $23$ solutions has $122$ elements.
The largest prime in the $17$ solutions of type (i) is $71$.
\cref{tab:solutions} lists, for each solution, its largest prime, its core size, and its elements outside the common part.

\begin{table}[ht]
\centering
\caption{The $23$ solutions of \cref{thm:B}. Each solution consists of the $23$ common elements listed in the text together with the $24$ elements in the last column. ``Core'' is the number of elements with both prime factors at most $73$.}
\label{tab:solutions}
\scriptsize
\setlength{\tabcolsep}{4pt}
\begin{tabularx}{\textwidth}{@{}rrrrX@{}}
\toprule
no. & max.\ prime & core & $P>73$ & remaining $24$ elements\\
\midrule
\tableinput{tables/solutions_table.tex}
\bottomrule
\end{tabularx}
\end{table}

\subsection{Correspondence between the claims and the programs}\label{sec:claimmatch}

\cref{claim:46,claim:47} are stated with exact rational quantities, and the programs implement them as follows.
\begin{itemize}
\item \textbf{Step~1.} By \cref{lem:step1}, every admissible core is output by the search. For the search written as a Lean function this is a theorem of the formalization (\cref{sec:lean}). That the C++ program computes the same function is not proved; it is supported by the comparisons of \cref{sec:leanrun}, where the compiled Lean function reproduces the C++ output.
\item \textbf{Shapes.} The shape enumeration generates every $(d,b)$ satisfying \ref{S1} and \ref{S2}. Since \ref{S4} is tested with a rounded-up right-hand side, every feasible shape passes.
\item \textbf{Candidate stars.} For every $A\subseteq W$ with $2\le\card{A}\le m$, the program removes all prime factors at most $73$ from $\nA(A)$ and factors the remainder completely. Each factor is certified prime by a deterministic test, and the product of the factors equals the remainder by construction. So every candidate star of \cref{def:candidate} is found. The residue condition for $d_q=1$ is tested exactly.
\item \textbf{Covers.} Covers are enumerated over index-increasing sequences of candidates with distinct centres, so every family of \cref{claim:47}(iii) is produced exactly once.
\item \textbf{Case $b=1$.} The program uses unordered pairs $\{A_1,A_2\}$, whereas \cref{claim:47}(iv) quantifies over ordered pairs. Swapping $(A_1,P_1)\leftrightarrow(A_2,P_2)$ replaces $X$ by $N/X$ and yields the same set.
\item \textbf{Magnitudes.} In \file{semiprime47.cpp} all intermediate integers in the $b=1$ branch are stored in $128$-bit words without explicit overflow checks. The Python recount (\cref{sec:python}) shows that in the actual budget-$47$ run the largest of them, $N$, has at most $59$ bits. Moreover, $b_1b_2$, $a_1b_2$ and $a_2b_1$ have at most $31$ bits and $t+a_1a_2$ at most $28$ bits. So no overflow occurred. In \file{semiprime48.cpp} all star and component arithmetic is $128$-bit with explicit overflow checks, and an overflow would be reported as unresolved.
\end{itemize}
This correspondence is an informal argument about the source code. It is \emph{not} machine-checked; see \cref{sec:lean}.

% =====================================================================
\section{Forty-eight terms}\label{sec:proof48}
% =====================================================================

With budget $\budget=48$ the same two steps determine all sets $T\subset\Pset$ with $\card{T}\le48$ and $\Rec{T}=1$.
Step~1 is the search of \cref{sec:proof46}, which is exhaustive by \cref{lem:step1}; in Lean this is \lean{core_mem_searchCores} with $\budget=48$.
Step~2 is that of \cref{sec:proof47}, with stars of at most $12$ leaves (\cref{lem:starsize}), and with one additional routine for shapes with $b=2$.

\subsection{Two large--large elements}\label{sec:twobb}

Suppose that $T$ satisfies \eqref{eq:setting} and that its shape has $b=2$.
Every large prime that divides no large--large element is the centre of a star, as in \cref{sec:component}.
The two large--large elements either share a prime, forming a path $P_0-P_1-P_2$, or they are disjoint, forming two edges $P_0P_1$ and $P_2P_3$.
For each of these large primes $P_i$ let $A_i=\Nb(P_i)\cap\Sset$ be its set of small neighbours, and let $a_i=\nA(A_i)$ and $b_i=\prod_{q\in A_i}q$ (with $a_i=0$, $b_i=1$ if $A_i=\varnothing$).
By \cref{cor:degree}, $A_i\ne\varnothing$ for every $P_i$ except the middle prime $P_1$ of a path.
The value of these elements is
\[
Z=1-\Rec{C}-\Rec{F}=\sum_i\frac{a_i}{b_iP_i}+\sum_{PP'}\frac1{PP'},
\]
where $F$ is the set of star elements and the last sum runs over the two large--large elements.
This is a sum of at most $5$ positive terms for a path and at most $6$ for two edges.
Hence one term is at least $Z/5$ or $Z/6$, respectively, and therefore some $P_i$ satisfies $P_i\le 6a_i/(b_iZ)$, or some large--large element $PP'$ satisfies $\min(P,P')\le\sqrt{6/Z}$.
For fixed $C$, attachment sets and stars, $Z$ is known, so at least one of the large primes is bounded by an explicit number $B$ (the program uses the sharper constant $5$ for paths).

The program tries every prime $p$ with $73<p\le B$ in every position:
\begin{itemize}
\item if $p$ is the middle of a path, then the two ends have the neighbours $A_0\cup\{p\}$ and $A_2\cup\{p\}$, so by \cref{lem:star} they are prime factors of $\nA(A_0\cup\{p\})=pa_0+b_0$ and of $pa_2+b_2$;
\item if $p$ is an end of a path, the other two primes form a single large--large edge in which $p$ is one more neighbour of the middle prime, and they are found through the divisor equation of \cref{prop:component};
\item if $p$ lies on one of two disjoint edges, its partner is a prime factor of $pa+b$ for the attachment data $(a,b)$ of the partner, and the other edge is found through the divisor equation.
\end{itemize}
A candidate is accepted only if its exact value is $Z$, and every accepted set is then verified exactly. Since in every completion some large prime is at most $B$, every completion with $b=2$ is found.
This argument is not formalized in Lean.

\subsection{The computation}

The run of \file{semiprime48.cpp} with budget $48$ is summarized in the last column of \cref{tab:results}.
Step~1 kept $1\,491\,334$ cores, $414$ of them with $U=\varnothing$.
Step~2 examined $116\,874\,859$ shapes; $601\,685$ of them, in $509\,055$ cores, satisfy \ref{S1}--\ref{S4}: $592\,627$ with $b=0$, $8\,985$ with $b=1$, $73$ with $b=2$ (in $73$ cores), and none with $b\ge3$.
For the $73$ shapes with $b=2$ the program solved $62\,169$ problems of the above form (choices of attachment sets and star covers); the largest bound $B$ was $1581$, and none of them yields a completion.
No core was unresolved and no probable prime was used; primality is decided by the Miller--Rabin test with the first $13$ prime bases, which is deterministic below $3.3\cdot10^{24}$ \cite{SorensonWebster2017}, and every number tested lies below that bound.
The run took $29$\,min\,$30$\,s with $4$ threads.

Step~2 finds every completion $G$ with $\card{G}\le48-\card{C}$, so the run finds every $T$ with $\card{T}\le48$ and $\Rec{T}=1$. It found $620$ sets with $48$ elements and, as a built-in consistency check, exactly the $23$ sets of \cref{thm:B}.

\begin{proof}[Proof of \cref{thm:C}, assuming the computation]
Let $\card{T}=48$ and $\Rec{T}\in\Z$; then $\Rec{T}=1$ by \cref{cor:sum-is-one}. The core of $T$ is $48$-admissible, so it is output by Step~1 (\cref{lem:step1}). Its shape is feasible (\cref{lem:shape}), and its large part is a completion of one of the kinds enumerated by Step~2: $b\le2$, since no feasible shape with $b\ge3$ occurs, and the completions with $b=0$, $1$ and $2$ are enumerated completely by \cref{sec:stars,sec:component,sec:twobb}. Hence $T$ is one of the $620$ listed sets. Conversely, each listed set has $48$ elements and reciprocal sum $1$ (verified exactly, see below).
The classification (i)--(iv) is read off from the list.
\end{proof}

\paragraph{Structure of the $620$ solutions.}
\cref{tab:solutions48} gives the classification of \cref{thm:C}; since no completion with $b=2$ exists, every solution has at most one large--large element.
The largest prime that occurs is $1\,210\,883$, in the set that contains
\[
11\cdot1\,210\,883,\quad 29\cdot1\,210\,883,\quad 37\cdot1\,210\,883,\quad 12\,109\cdot1\,210\,883,\quad 17\cdot12\,109 ;
\]
its largest element $12\,109\cdot1\,210\,883=14\,662\,582\,247$ is the largest element of any of the $620$ sets.
Exactly $12$ numbers lie in every solution: $6$, $10$, $14$, $15$, $21$, $22$, $26$, $33$, $34$, $38$, $39$, $46$.
In particular, the $48$-term representation of $1$ attributed to Johnson \cite{Johnson1978} in \cite{Watanabe2020} must be one of the $620$ sets. We have not been able to consult \cite{Johnson1978} and have not identified it.

\begin{table}[ht]
\centering
\caption{The $620$ sets $T\subset\Pset$ with $\card{T}=48$ and $\Rec{T}=1$, by the primes larger than $73$ that they contain (\protect\file{solutions48.txt}; recomputed by \protect\file{paper/verification/verify48_facts.py}).}
\label{tab:solutions48}
\small
\begin{tabular}{@{}lr@{}}
\toprule
primes larger than $73$ & number of sets\\
\midrule
none & $397$\\
one prime $P$, the centre of a star & $187$\\
two primes, both centres of stars & $18$\\
two primes $P_1,P_2$ with $P_1P_2\in T$ & $18$\\
\midrule
total & $620$\\
\bottomrule
\end{tabular}
\end{table}

\paragraph{Verification.}
Every listed set was verified in three ways: by \file{semiprime48.cpp} with big-integer fractions; by \file{verify48.py}, with its own factorization, a deterministic primality test and Python fractions; and by \file{paper/verification/verify48_facts.py}, which uses trial division, checks $\Rec{T}=1$ with Python fractions, checks the congruences \eqref{eq:local} at every vertex, and recomputes all numbers of this subsection.
The Python program \file{complete48.py} is an independent implementation of Step~2 for budget $48$, including its own code for $b=2$. Run on all $1\,491\,334$ cores of the dump \file{cores48.txt.gz}, it finds the same $620+23$ sets and no unresolved core.
Both implementations of the routine for $b=2$ recover planted examples (a path and a pair of disjoint edges) with identical sets of completions.
An earlier run, made before the $b=2$ routine existed, reported exactly the $73$ cores with $b=2$ as unresolved and produced a core dump identical to that of the final run.

\paragraph{Status.}
\cref{thm:C} is computer-assisted in the same sense as \cref{thm:A,thm:B}, with two differences: its Step~2, including the routine for $b=2$, is not stated in Lean, and the $620$ sets are verified exactly by the C++ and Python programs but not in Lean. The Lean formalization covers the reduction to Step~2 for budget $48$: the integrality of $\Rec{T}$ implies $\Rec{T}=1$ (\lean{recipSum_eq_one_of_integral_48}), and the core of every solution is in the output of the verified search (\lean{core_mem_searchCores}).

% =====================================================================
\section{Exact and independent verification}\label{sec:verification}
% =====================================================================

We distinguish five layers of evidence:
\begin{enumerate*}[label=(\arabic*)]
\item the optimized C++ searches, in two generations (without and with the loss bound);
\item independent Python re-implementations of Step~2;
\item exact verification of the solutions;
\item the search written as a Lean function, proved exhaustive in Lean and executed for comparison;
\item the Lean formalization of the reduction (\cref{sec:lean}).
\end{enumerate*}
We also report checks performed specifically for this paper. The verification of the $620$ sets of \cref{thm:C} is described in \cref{sec:proof48}.

\subsection{Exact verification of the 23 solutions}\label{sec:exact}

Each of the $23$ sets was verified in five independent ways.
\begin{itemize}
\item \file{semiprime47.cpp} and \file{semiprime48.cpp} factor every element by trial division, certify both factors with a deterministic primality test, and compute the reciprocal sum as a big-integer fraction over the product of all primes involved.
\item \file{complete47.py} checks the sum with Python's \code{fractions.Fraction}; so does \file{verify48.py}, which also checks that the $47$-element sets found by the budget-$48$ run are exactly these $23$.
\item In Lean, \lean{Sol23_valid} proves that each set consists of $47$ squarefree semiprimes, with explicit factorizations and primality proved by \code{norm\_num}, and has reciprocal sum exactly $1$. \lean{Sol23_card} proves that the $23$ sets are pairwise distinct.
\item For this paper, the script \file{paper/verification/verify_paper_facts.py} rechecks everything with exact integer arithmetic and trial division. It checks distinctness of the elements and of the sets, the squarefree-semiprime property, $\Rec{T}=1$, and the congruences \eqref{eq:local} at every vertex. It also recomputes all numbers quoted in \cref{sec:prelim,sec:stars} and generates all tables of this paper from the data.
\end{itemize}

\subsection{Independent implementation of Step 2}\label{sec:python}

The Python program \file{complete47.py} is a separate implementation of Step~2. It uses exact rationals (\code{fractions.Fraction}) and unbounded integers throughout, its own primality test and factorization (trial division followed by Pollard's rho method \cite{Pollard1975}), and its own cover and divisor enumeration.
It imposes no bound on the number of leaves of a candidate star, so it examines a superset of the candidates. It also recomputes $U$ from each core and compares it with the value in the dump.
Run on the file \file{cores47.txt.gz} of all $22\,382$ cores kept by Step~1 with $\budget=47$, it finds the same $23$ solutions and no unresolved core (\file{run47_python_crosscheck.txt}).
For this paper we reran it with additional counters (\file{paper/verification/analyze47.py}, identical logic). It reproduces every count in the budget-$47$ column of \cref{tab:results}: $1\,107\,702$ shapes, of which $8\,049$ survive in $7\,432$ cores, and $125$ candidate stars (with or without the bound of $11$ leaves). It also gives the split by $b$ and the magnitudes quoted in \cref{sec:claimmatch}. The run takes about half a minute.

The Python program \file{complete48.py} does the same for budget $48$, with its own implementation of the routine of \cref{sec:twobb} (\cref{sec:proof48}).

\emph{Limitations.} The Python programs check Step~2 only. They consume the output of the C++ Step~1 and cannot detect a core that Step~1 failed to output; that gap is addressed by \cref{lem:step1}, its formal version, and the comparisons of \cref{sec:leanrun}. They were written within the same project and implement the same mathematical reduction (\cref{sec:core,sec:stars}), so they are independent \emph{implementations}, not independent \emph{arguments}.

\subsection{Re-executions and the budget-46 cross-check}\label{sec:rerun}

%% generated by verification/make_tables.py from verification/rerun/*.log -- do not edit by hand
\paragraph{Re-execution.} For this paper the budget-$46$ program \file{semiprime_reciprocals.cpp} was compiled and executed again with $4$ threads (GCC~13.3.0; a cloud container with $4$ virtual CPUs, Intel Xeon 2.10\,GHz). The output is identical to the recorded \file{proof_run_output.txt}, including the node count $139\,830\,520\,211$, except for the elapsed time ($1531$\,s, with other jobs running concurrently); log: \file{paper/verification/rerun/}.

\paragraph{Budget-46 cores and a second Step~2.} The engine of \file{semiprime47.cpp}, included unmodified in a small driver, was run with budget $46$; it writes every core kept by Step~1 to a file. It visited $139\,830\,520\,211$ nodes, kept $178$ cores, examined $3\,777$ shapes of which $61$ survive, found $0$ candidate stars, $0$ unresolved cores and $0$ solutions, in agreement with \cref{tab:results}. The Python Step~2 (\file{analyze47.py}, the logic of \file{complete47.py}) applied to these $178$ cores with budget $46$ likewise finds $3\,777$ shapes satisfying \ref{S1}--\ref{S3}, $61$ satisfying \ref{S1}--\ref{S4} (all with $b=0$), $0$ candidate stars (also without the bound of $10$ leaves), $0$ unresolved shapes and $0$ solutions. This archives the Python check of the budget-$46$ Step~2 mentioned in the repository documentation. Since the Step~1 code of both programs is the same, this is a second \emph{execution} of Step~1, not an independent implementation.

\paragraph{Re-execution of the budget-47 search.} The program \file{semiprime47.cpp} was also executed again, with $4$ threads and the option \file{--dump}, on the same machine. All printed quantities coincide with the recorded \file{run47_output.txt}: the node count $1\,795\,181\,713\,099$, the $22\,382$ cores, the shape and candidate counts, $0$ unresolved cores and $0$ probable primes. The list of $23$ solutions is identical, and the new core dump coincides line by line (as a set of $22\,382$ lines) with \file{cores47.txt.gz}. Only the elapsed time differs ($19\,388$\,s for the search). Files: \file{paper/verification/rerun47/}.


\paragraph{Old and new search.} By \cref{lem:sameoutput} the searches with and without \ref{P4} have the same output. This was checked directly (\file{validation48/}): \file{semiprime48.cpp} reproduces the archived core dumps of the earlier programs line for line, $178$ cores for budget $46$ (with $2.6\cdot10^{7}$ instead of $1.4\cdot10^{11}$ nodes) and $22\,382$ cores for budget $47$ (with $6.0\cdot10^{8}$ instead of $1.8\cdot10^{12}$ nodes). For budgets $40$--$45$ on the real instance both programs keep no core, and on $15$ reduced instances ($\Sset$ = the primes up to $13$, $17$, $19$ and $23$, with up to $21$ million cores) the two programs produce identical dumps. On these small instances \ref{P4} prunes nothing, so they test the bookkeeping of the new code rather than the pruning itself.

\subsection{Validation of Step 1 by brute force on reduced instances}\label{sec:bruteforce}

Step~1 has no second full-scale implementation in C++ or Python (the two C++ generations share most of their code, and the Lean function of \cref{sec:leanrun} is too slow for budgets above $46$). To test the C++ code, we ran the engine of \file{semiprime47.cpp} on reduced instances, with $\Sset$ replaced by the primes up to $13$, $17$ and $19$ and with various budgets. We compared its output with a brute-force enumeration of \emph{all} $2^{15}$, $2^{21}$ and $2^{28}$ cores, respectively.
For this test only, a copy of the program with Step~2 disabled was used; the Step~1 code is unchanged.
For every core the brute force computes two sets, using the same fixed-point rounding as the program:
\begin{itemize}
\item $\mathrm{SEG}$: the cores satisfying \ref{P1}, $\Rec{C}\le1$ exactly, and the rounded test \eqref{eq:leaftest} for \emph{every} upper segment of $U$ (\cref{def:admissible});
\item $\mathrm{FIN}$: the cores satisfying \ref{P1}, \ref{P2} and the final test \eqref{eq:leaftest}, i.e.\ only the case $U'=U$.
\end{itemize}
\cref{lem:step1} and the form of the leaf test predict $\mathrm{SEG}\subseteq D\subseteq\mathrm{FIN}$ for the set $D$ of cores output by the search.
This was confirmed in all $48$ configurations tested: $12$ pairs of cutoff and budget, the default table size of $8$ low primes and a reduced size of $3$ (which exercises the explicit recursion), and $1$ and $3$ threads. The search output never contained a duplicate.
A negative control, which compared a budget-$45$ search with the reference sets for budget $50$, was correctly reported as a failure.
Representative sizes are in \cref{tab:bruteforce}. They also show that $\mathrm{FIN}$ can be much larger than $\mathrm{SEG}$, which quantifies \cref{rem:segments}.

\begin{table}[ht]
\centering
\caption{Brute-force validation of Step 1 on reduced instances ($\Sset$ = primes up to $s$). $D$ is the output of the search, with the default table of the $8$ smallest primes and with a table of $3$ primes. The output did not depend on the number of threads. In all cases $\mathrm{SEG}\subseteq D\subseteq\mathrm{FIN}$.}
\label{tab:bruteforce}
\small
\begin{tabular}{@{}rrrrrrr@{}}
\toprule
 & & & & \multicolumn{2}{c}{$\card{D}$, table size} & \\
\cmidrule(lr){5-6}
$s$ & edges & budget & $\card{\mathrm{SEG}}$ & $8$ & $3$ & $\card{\mathrm{FIN}}$\\
\midrule
\tableinput{tables/bruteforce_table.tex}
\bottomrule
\end{tabular}
\end{table}

These tests exercise the same code path as the real search, including the table-driven enumeration, the residue bookkeeping, the $U$-accounting, all pruning tests and the parallel work distribution, on instances small enough to enumerate completely. They do not replace an independent full-scale implementation of Step~1.
The same brute-force check was passed by \file{semiprime48.cpp} on the reduced instances with $\Sset$ = the primes up to $13$, $17$ and $19$ (\file{validation48/step1_equivalence_output.txt}).

\subsection{The search as a Lean function}\label{sec:leanrun}

The Lean function \lean{Inst.search} of \cref{sec:lean} follows Step~1 of \file{semiprime48.cpp}: the same order of the primes, the same splitting into explicit recursion and a final step for the $8$ smallest primes, the same tests \ref{P1}--\ref{P4} with the same integer tables, and the same final test. It is written for clarity, not speed: cores are finite sets of index pairs, residues are recomputed at every node, and the final step runs through all $256$ subsets of the $8$ smallest primes instead of using a table indexed by residues.
Lean proves that its output contains every admissible core (\cref{sec:lean}). Its output is not certified to be \emph{equal} to that of the C++ program. To test the correspondence we compiled the Lean function to native code and compared its output with that of \file{semiprime48.cpp} (option \code{-{}-dump}; harness \file{validation48/dfs48_harness.cpp}):
\begin{itemize}
\item on $12$ reduced instances ($\Sset$ = the primes up to $13$ with budgets $40$, $45$, $50$, $60$; up to $17$ and up to $19$ with budgets $40$, $45$, $50$; up to $23$ with budgets $40$, $45$; instance generator \lean{mkInst}), the two outputs coincide as sets of cores; the largest outputs have $274\,486$ cores ($\Sset$ = primes up to $23$, budget $45$) and $478\,084$ cores ($\Sset$ = primes up to $19$, budget $50$);
\item on the real instance ($\Sset$ = the primes up to $73$), for budgets $40$ to $45$ both keep no core, and for budget $46$ both keep the same $178$ cores.
\end{itemize}
The compiled Lean function is several hundred to several thousand times slower than the C++ program (for budget $46$: about $8.5$ CPU hours, against $4.6$\,s for the C++ search with one thread), mainly because it recomputes all residues at every node and evaluates all $256$ choices at every final step. For budget $46$ it was therefore run in several processes: \lean{search_eq_rounds} (\file{SearchSplit.lean}) proves that the search is the concatenation, in order, of the searches below the $890\,624$ work items obtained by $26$ rounds of refinement of the root; the heaviest item was refined further in the same way (\lean{flatMap_expandItem}), and every piece was evaluated by exactly one process (\file{validation48/lean_runner/README.md}). Budgets $47$ and $48$ are out of reach for the Lean function.
The comparison script and its output are \file{validation48/lean_vs_cpp.sh} and \file{validation48/lean_vs_cpp_output.txt}.

\subsection{Self-tests}

The recorded self-tests (\file{selftest_output.txt}) provide further consistency checks.
\begin{itemize}
\item The engine, run without unsatisfied primes and with budget $47$, finds exactly $17$ solutions. These are the $17$ solutions of \cref{thm:B}(i).
\item The completion routine reconstructs a planted solution with the cutoff lowered to $53$ and to $43$. In the second case the shape has positive excess.
\item The full pipeline with cutoff $53$ and budget $47$ finds $21$ solutions, all among the $23$. With cutoff $53$ it also leaves $2902$ unresolved cores; this is expected, because the parity and value tests are much weaker there, and it is not an error.
\end{itemize}
The engine of \file{semiprime47.cpp} also recovers a planted component with $b=1$ through the divisor equation (\file{run47_output.txt}).

\subsection{What is not independently verified}

The following are \emph{not} independently verified:
\begin{itemize}
\item the full-scale Step~1 for budgets $47$ and $48$, which was executed only by the C++ programs; the Lean function whose exhaustiveness is proved was executed for budgets $40$ to $46$ only (\cref{sec:leanrun});
\item the correctness of the compilers and the hardware;
\item the informal correspondence of \cref{sec:claimmatch} between \cref{claim:46,claim:47} and the source code, and the corresponding correspondence for Step~2 with budget $48$.
\end{itemize}
The budget-$46$ Step~1 was, however, executed five times with identical results: in the recorded run, in a re-execution, through the engine of \file{semiprime47.cpp}, by \file{semiprime48.cpp}, and as the compiled Lean function. The budget-$47$ search was executed three times (twice without and once with \ref{P4}), and all runs produced identical core dumps and identical solution lists (\cref{sec:rerun}). The budget-$48$ search was executed twice, with identical core dumps (\cref{sec:proof48}).

% =====================================================================
\section{Lean formalization}\label{sec:lean}
% =====================================================================

The directory \file{lean/} contains a Lean~4 \cite{Lean4} development based on Mathlib \cite{Mathlib2020}. It uses the toolchain \file{leanprover/lean4:v4.35.0-rc3} and the Mathlib commit
\begin{center}\ttfamily 3cb72cfd416d1b5ec4b930d67648ef036a5df24f\end{center}
and consists of $18$ files with about $5{,}100$ lines and about $320$ theorems and lemmas: nine files for the analytic reduction and the solutions, and nine files for the search (\file{TopSum}, \file{Search}, \file{LossBound}, \file{SearchCorrect}, \file{SearchTables}, \file{SearchTablesProof}, \file{SearchReal}, \file{MainSearch}, \file{SearchSplit}).
It contains no \code{sorry}, \code{admit}, \code{axiom} or \code{native\_decide}. The main theorems depend only on the axioms \code{propext}, \code{Classical.choice} and \code{Quot.sound}, as reported by \code{\#print axioms} (\file{lean/AxiomsCheck.lean}; output in \file{paper/verification/lean_axioms_output.txt}).
Auto-bound implicit variables are disabled.
The development builds with \code{lake build}.

\paragraph{Formally verified: the analytic reduction.}
The following results are proved in Lean without any computational hypothesis:
\begin{itemize}
\item the local criterion (\cref{lem:local}; \lean{integral_iff_localCond});
\item the exchange lemma, the list of the $47$ smallest elements of $\Pset$ together with the fact that it is complete, the inequalities of \cref{prop:Hk}, and \cref{cor:sum-is-one} (\lean{recipSum_le_Aupto}, \lean{sp_le_158}, \lean{H37_lt_one}, \lean{H38_gt_one}, \lean{H47_lt_two}, \lean{recipSum_eq_one_of_integral}, \lean{recipSum_eq_one_of_integral_48}, \lean{card_ge_38});
\item \cref{lem:U} (\lean{card_Uset_le}), \cref{lem:value} in the general form for every $U'\subseteq U$ (\lean{value_lower_subset}), and all parts of \cref{lem:shape} (\lean{dcount_pos_of_mem_U}, \lean{dcount_ne_one_of_not_mem_U}, \lean{excess_le}, \lean{parity_two}, \lean{shape_value_bound});
\item the star lemmas and the star-size lemma (\lean{star_dvd}, \lean{star_card_ne_one}, \lean{star_size}), the decomposition of $G$ into stars and large--large edges (\lean{bigPart_eq}, \lean{recipSum_bigPart_eq}, \lean{dcount_eq_mult}), and the finiteness of the candidate stars (\lean{starCand_finite});
\item all parts of \cref{prop:component} (\lean{comp_value}, \lean{comp_local_iff}, \lean{comp_divisor_eq}, \lean{comp_enum}, \lean{comp_of_divisor}, \lean{comp_pairs_finite});
\item the statement that the core and shape of every solution pass the tests of the search (\lean{admissible_of_real}, \lean{shapeOK_of_real}).
\end{itemize}

\paragraph{Unconditional verification of the solutions.}
\lean{Sol23_valid} proves that each of the $23$ sets consists of $47$ squarefree semiprimes and has reciprocal sum $1$, and \lean{Sol23_card} proves that there are $23$ distinct sets.
The numerical facts are checked by the kernel through \code{decide} and \code{norm\_num}.
In particular, the existence part of \cref{thm:A} (sharpness) is fully formally verified.

\paragraph{The search, and the proof that it is exhaustive.}
Step~1 is written as a Lean function, and its exhaustiveness is a theorem.
\begin{itemize}
\item \file{Search.lean} defines an \emph{instance} (the primes of $\Sset$, the budget, the least large prime, the list $v_1,v_2,\dots$ of \cref{sec:step1}, $\theta$, the share numerators and the number $8$ of primes of the final step) and the search \lean{Inst.search} as a recursive function over the levels $(j,i)$ of \cref{alg:step1}. A core is a finite set of index pairs $(a,b)$, $a<b$, standing for $p_ap_b$. The tests \ref{P1}--\ref{P4} and the final test \eqref{eq:leaftest} are Boolean functions with exact natural-number and integer arithmetic. The expensive tables ($\mathrm{VB}$, the maxima $M$ of \ref{P3}, $V_\theta$ and the local bounds of \ref{P4}) are an argument of the search.
\item \file{SearchTables.lean} computes these tables with the formulas of the C++ program, the local bounds by a knapsack over residues (\lean{dpRun}).
\item \lean{mem_search} (\file{SearchCorrect.lean}): for every instance and every family of tables with the soundness properties \lean{TablesOK}, the output of the search contains every core that satisfies the rounded admissibility condition \lean{IntAdm} (the budget, $\lo(C)\le2^{96}$, and the rounded value test \eqref{eq:leaftest} for every upper segment of $U$). The proof follows the path to the core in the search tree and shows that no test fails there; for \ref{P4} it uses \lean{loss_bound} (\file{LossBound.lean}), the formal version of \cref{prop:loss}, including the rounding of the shares.
\item \lean{mkTables_ok} (\file{SearchTablesProof.lean}): the computed tables satisfy \lean{TablesOK}; in particular the knapsack computes upper bounds of the local maxima.
\item \lean{intAdm_of_admissible} (\file{SearchReal.lean}): for $\Sset$ = the primes up to $73$ and $\budget\le48$ (the instance \lean{inst73}), every core with \code{Admissible K C} (\cref{def:admissible}, with exact rationals) satisfies \lean{IntAdm}. This includes the facts that the residues computed by the search are the residues $\rho$ of the core, that $\lo(n)\le2^{96}/n\le\up(n)$, and that the $49$ listed numbers $158,166,\dots,471$ are the $49$ smallest elements of $\Pset$ with a large prime factor, so that $\mathrm{VB}(x)\ge2^{96}V(x)$.
\item Hence \lean{core_mem_search}: every admissible core (as its edge set \code{edgesOf C}) is in the output of \code{(inst73 K).search (inst73 K).mkTables}, for every $\budget\le48$; this is \cref{lem:step1}. With \code{searchCores K} the list of cores of that output, \lean{core_mem_searchCores} states that the core of every nonempty $T\subset\Pset$ with $\card{T}\le\budget\le48$ and $\Rec{T}\in\Z$ lies in \code{searchCores K}.
\end{itemize}
The search is thus a Lean function whose output is proved to contain the core of every solution; its output is not computed inside Lean's kernel. The function is also executable, and the compiled function reproduces the output of the C++ program for budgets $40$ to $46$ and on reduced instances (\cref{sec:leanrun}).

\paragraph{Conditional completeness.}
Step~2 is not reproduced in Lean. Instead, \cref{claim:46,claim:47} are stated as the propositions \lean{Search46} and \lean{Search47} (reproduced in \cref{app:lean}). These are ordinary definitions used as hypotheses, not axioms. Their weaker versions \lean{Search46'} and \lean{Search47'} quantify only over the cores in \code{searchCores 46} and \code{searchCores 47}, the output of the verified search. By the exhaustiveness theorem, \lean{Search46'} implies \lean{Search46} and \lean{Search47'} implies \lean{Search47} (\lean{search46_of}, \lean{search47_of}). Lean proves:
\begin{itemize}
\item \lean{no_integral_sum_le_46}, which is \cref{thm:A} with \lean{Search46} as a hypothesis, its reformulation \lean{card_ge_47_of_integral}, and the same statements with the hypothesis \lean{Search46'} (\lean{no_integral_sum_le_46'}, \lean{card_ge_47_of_integral'});
\item \lean{solutions47_iff}, which is \cref{thm:B} with \lean{Search47} as a hypothesis; \lean{solutions47_integral_iff} states the version for integral sums, and \lean{solutions47_ncard} states that the set of all $47$-term representations of $1$ has exactly $23$ elements; the primed versions use the hypothesis \lean{Search47'}.
\end{itemize}
Consequently the complete proofs of \cref{thm:A,thm:B} are formally verified \emph{modulo} \lean{Search46'} and \lean{Search47'}: the enumeration of the cores is no longer an assumption, and what remains assumed is the outcome of Step~2 on the cores output by the verified search. These hypotheses are established by the computations of \cref{sec:proof46,sec:proof47}. The following points are not verified in Lean:
\begin{itemize}
\item the fact that the C++ Step~1 outputs (at least) the cores of \code{searchCores K}; for $\budget=46$ this was checked by running the compiled Lean function, and for $\budget=47$ it rests on the C++ program implementing the same function (\cref{sec:leanrun});
\item the fact that the C++ Step~2 implements \cref{claim:46,claim:47} correctly (\cref{sec:claimmatch});
\item the fact that the recorded runs produced the stated output.
\end{itemize}
For \cref{thm:C}, Lean provides the reduction to the cores of \code{searchCores 48} (\lean{core_mem_searchCores}), but the Step~2 statement for budget $48$ (in particular the case $b=2$ of \cref{sec:twobb}) is not formalized, and the $620$ sets are not verified in Lean.

While the first part of the development was being written, the formalization exposed the two corrections of \cref{rem:segments,rem:component}. The Step~1 hypothesis is stated with the upper-segment condition of \cref{def:admissible}. With only the final value test, it would not describe what the search checks.

% =====================================================================
\section{Computational reproducibility}\label{sec:repro}
% =====================================================================

\paragraph{Arithmetic.}
None of the C++ programs uses a floating-point type (\code{float}, \code{double} or \code{long double}), nor any floating-point library function.
\begin{itemize}
\item Congruences are computed with machine integers modulo primes at most $73$.
\item Exact rational comparisons (the values $H_k$, the verification of solutions, the exact value tests of covers) use either a small arbitrary-precision natural-number type or $128$-bit integers over the common denominator $D=\prod_{p\in\Sset}p\approx4.0\cdot10^{28}$.
\item Bounds on sums of reciprocals use the $96$-bit fixed-point numbers $\up$ and $\lo$ of \cref{sec:step1}. Upper bounds are always formed from $\up$ and lower bounds from $\lo$. Every pruning decision compares an upper bound of a quantity that must be at least $1$ with $2^{96}$, or a lower bound of a quantity that must be at most $1$ with $2^{96}$.
\item The loss bound \ref{P4} uses signed $128$-bit integers; the shares are rounded up and the knapsack tables are exact maxima, so it is an upper bound as in \cref{prop:loss}.
\end{itemize}
Rounding can therefore only make a test \emph{weaker}. It may keep a core or shape that exact arithmetic would discard, and this costs only time; it can never discard an admissible core or a feasible shape. All fixed-point sums are below $3\cdot2^{96}<2^{128}$.
The Python program uses exact rationals throughout.

\paragraph{Primality and factorization.}
The earlier programs decide primality of $64$-bit integers with the Miller--Rabin test \cite{Rabin1980} to the bases $2,3,5,\dots,37$ (the first $12$ primes).
This test is deterministic for all $n<\psi_{12}=318\,665\,857\,834\,031\,151\,167\,461\approx3.2\cdot10^{23}$ \cite{SorensonWebster2017}, and in particular for all $n<2^{64}$.
The $b=1$ branch of \file{semiprime47.cpp} contains a probabilistic test for integers of at least $2^{64}$. It was never used: the run reports $0$ probable primes.
\file{semiprime48.cpp} uses the first $13$ prime bases, which give a deterministic test below $\psi_{13}\approx3.3\cdot10^{24}$ \cite{SorensonWebster2017}, and counts every number above that bound that passes the test as a probable prime; all its runs report $0$ probable primes.
Factorizations use trial division and Pollard's rho method \cite{Pollard1975}. The correctness of the result does not depend on the heuristics of the method, because every returned factor is certified prime and the factors multiply to the input by construction; only termination does, and the runs terminated.
The inverses modulo $p$ are computed as $a^{p-2}\bmod p$.

\paragraph{Build and run.}
The programs need a C++11 compiler with \code{unsigned \_\_int128} (GCC or Clang):
{\small
\begin{verbatim}
g++ -O3 -march=native -pthread -o semiprime48 semiprime48.cpp
./semiprime48 4 --budget 46             # Theorem A, 4 threads (~2 s)
./semiprime48 4 --budget 47             # Theorem B (~1 min)
./semiprime48 4 --selftest --dump cores48.txt   # Theorem C (~30 min)
python3 verify48.py run48_output.txt
gunzip -k cores48.txt.gz && python3 complete48.py cores48.txt 48 4
bash validation48/step1_equivalence.sh 4        # old vs new Step 1
# earlier programs:
g++ -O3 -march=native -pthread -o semiprime_reciprocals semiprime_reciprocals.cpp
./semiprime_reciprocals 4               # budget 46
g++ -O3 -march=native -pthread -o semiprime47 semiprime47.cpp
./semiprime47 4 --selftest --dump cores47.txt
gunzip -k cores47.txt.gz && python3 complete47.py cores47.txt 47
\end{verbatim}
}
The recorded runs of \file{semiprime48.cpp} used $4$ threads and took $2.1$\,s (budget $46$), $65$\,s (budget $47$) and $1770$\,s (budget $48$, 29\,min\,30\,s) for search and completion.
The recorded runs of the earlier programs also used $4$ threads. The budget-$46$ search took $1179$\,s (19\,min\,40\,s wall-clock), and the budget-$47$ run took $20\,898$\,s for the search and 5\,h\,48\,min wall-clock in total.
The Lean search is compiled to native code with the scripts in \file{validation48/lean_runner/}, and \file{validation48/lean_vs_cpp.sh} compares its output with that of \file{semiprime48.cpp}.
The resident memory of the search is only a few megabytes; the virtual size of about $300$\,MB reported in the repository documentation consists mostly of thread stacks.
The recorded outputs do not state the processor model.
The re-executions of \cref{sec:rerun} used a cloud container with $4$ virtual CPUs (Intel Xeon, 2.10\,GHz) and GCC~13.3.0. There the budget-$47$ search took $19\,388$\,s with $4$ threads.

\paragraph{Files.}
The repository (\repourl) contains the following files.
\begin{itemize}
\item Programs: \file{semiprime48.cpp}, \file{complete48.py} and \file{verify48.py}; the earlier programs \file{semiprime_reciprocals.cpp}, \file{semiprime47.cpp} and \file{complete47.py}.
\item Recorded outputs:
  \begin{itemize}
  \item of \file{semiprime48.cpp}: \file{run48_output.txt}, \file{validation48/run46_output.txt} and \file{validation48/run47_output.txt};
  \item of the Python checks: \file{run48_python_crosscheck.txt}, \file{verify48_output.txt} and \file{run47_python_crosscheck.txt};
  \item of the earlier programs: \file{proof_run_output.txt}, \file{selftest_output.txt} and \file{run47_output.txt}.
  \end{itemize}
\item Data: the core dumps \file{cores47.txt.gz} and \file{cores48.txt.gz}, and the lists \file{solutions47.txt} and \file{solutions48.txt}.
\item Validation of the new search: \file{validation48/} (old versus new Step~1, brute force, the case $b=2$, Lean versus C++).
\item The Lean development \file{lean/}.
\item The files of this paper in \file{paper/}, including the verification scripts, the brute-force test and the logs of the re-execution (\file{paper/verification/}).
\end{itemize}

% =====================================================================
\section{Discussion}\label{sec:discussion}
% =====================================================================

We have shown that at least $47$ distinct squarefree semiprimes are needed to obtain an integral sum of reciprocals, and that exactly $23$ sets of $47$ such numbers have reciprocal sum $1$. With the same method we found that exactly $620$ sets of $48$ such numbers have reciprocal sum $1$.
The minimality result confirms Watanabe's conjecture \cite{Watanabe2020}.

The main mathematical point is the finite reduction. The local criterion (\cref{lem:local}) turns integrality into a condition at each vertex of a graph. The core decomposition bounds the number of unsatisfied small primes. The star lemma and the component equation bound the large primes by divisibility, not by size.
That large primes occur, for example $3779$, shows that a bound of this kind is unavoidable.
A single large--large edge is handled by an explicit factorization identity, and two such edges by a size argument that bounds one of the large primes (\cref{sec:twobb}). Configurations with three or more such edges are excluded here only by the computation: no feasible shape with $b\ge3$ exists for budget $48$. For larger budgets this case would need further ideas.
On the computational side, the loss bound (\cref{sec:lossbound}) makes Step~1 thousands of times faster. It combines a Lagrangian relaxation of the budget constraint with the congruences that every prime not yet processed must eventually satisfy, and it leaves the output of the search unchanged.

\paragraph{Limitations.}
The proofs are computer-assisted.
\begin{itemize}
\item Step~1 is now a Lean function with a formal proof of exhaustiveness. The remaining gap is that the full-scale executions for budgets $47$ and $48$ were made by the C++ program, which is meant to compute the same function but is not verified; on every instance on which both were run, including the real instance with budgets $40$ to $46$, the two produce the same cores. The two generations of C++ programs share most of their Step~1 code.
\item Step~2 has been implemented twice, in C++ and in Python, and the implementations agree for all three budgets.
\item The reduction, the exhaustiveness of Step~1 and the $23$ solutions are verified in Lean. The connection between the Lean hypotheses \lean{Search46'} and \lean{Search47'} and the C++ Step~2 rests on the informal argument of \cref{sec:claimmatch}. For budget $48$, Step~2 (including the case $b=2$) and the $620$ solutions are not formalized.
\end{itemize}
A fully formal proof would require the computation to emit a certificate that can be checked in Lean, such as the list of cores output by Step~1 together with a refutation or completion of each feasible shape; alternatively, Step~2 could itself be written and verified as a Lean function, as was done here for Step~1, and both executed in a kernel-checkable way. This is a natural direction for future work.
An independent reimplementation of the search, preferably by other authors and with a different enumeration strategy, would also strengthen the result considerably.

\paragraph{Further questions.}
The same method applies to other targets and budgets. Examples include the minimal number of terms for the integer $2$, or for denominators with exactly three distinct prime factors (compare \cite{ButlerErdosGraham2015}). With the loss bound the cost of Step~1 still grows by a factor of about $20$ per unit of budget: $2.6\cdot10^{7}$, $6.0\cdot10^{8}$ and $1.3\cdot10^{10}$ nodes for budgets $46$, $47$ and $48$.

% =====================================================================
\section{Disclosure of the use of generative AI}\label{sec:ai}
% =====================================================================

% This statement is intended to be easily editable to follow the policy of the target journal.
Generative AI tools, specifically Claude Code (Anthropic), an AI coding assistant based on large language models, were used substantially throughout this project, and not only for preparing the manuscript.
Under the direction of the authors, the AI tool developed the mathematical reduction presented here, including the loss bound of \cref{sec:lossbound} and the treatment of two large--large elements in \cref{sec:twobb}, and it wrote the following:
\begin{itemize}
\item the C++ programs and the Python cross-checks;
\item the Lean formalization, including the search as a Lean function and the proof of its exhaustiveness;
\item the documentation files of the repository;
\item the verification scripts written for this paper;
\item a first draft of this manuscript and of its revision.
\end{itemize}
The exploratory mathematical reasoning was also carried out with the tool.
The two corrections of \cref{rem:segments,rem:component} were found during the AI-assisted formalization and a later audit of the code.

All mathematical claims in this article were checked by the verification procedures described in \cref{sec:proof48,sec:verification,sec:lean}:
\begin{itemize}
\item exact computation;
\item second implementations of Step~2;
\item brute-force validation of Step~1 on reduced instances, and the comparison of the two generations of the search;
\item a Lean~4 formalization of the analytic reduction, of the exhaustiveness of Step~1 and of the $23$ solutions, and the execution of the verified Lean search for budgets $40$ to $46$.
\end{itemize}
These checks were themselves designed and executed with the assistance of the same tool.
The extent of the computer-assisted part and of the formally verified part, and the exact status of every result, are stated explicitly in \cref{fig:status} and \cref{sec:proof46,sec:proof47,sec:proof48,sec:verification,sec:lean}.
[\emph{Editable by the authors: describe here any independent human checking of the code, the computations and the proofs that has been carried out.}]
No AI system is an author of this work. The authors take full responsibility for its content.

% =====================================================================
\section*{Acknowledgements}
% =====================================================================
% Intentionally left blank -- to be completed by the authors.

\bibliographystyle{plainnat}
\bibliography{references}

% =====================================================================
\appendix
\section{The 23 representations of 1 with 47 terms}\label{app:solutions}
% =====================================================================

The $23$ sets $T\subset\Pset$ with $\card{T}=47$ and $\Rec{T}=1$, in the order of \file{solutions47.txt}. Solutions $1$, $2$, $3$, $5$, $11$ and $12$ contain a prime larger than $73$ (\cref{tab:bigprimes}).

\medskip
% generated by verification/verify_paper_facts.py -- do not edit by hand
\solitem{1}{6, 10, 14, 15, 21, 22, 26, 33, 34, 35, 38, 39, 46, 51, 55, 57, 58, 62, 65, 69, 74, 77, 82, 85, 87, 91, 93, 95, 111, 115, 119, 123, 133, 143, 145, 155, 203, 219, 221, 287, 299, 391, 481, 1299, 2117, 16021, 31609}
\solitem{2}{6, 10, 14, 15, 21, 22, 26, 33, 34, 35, 38, 39, 46, 51, 55, 57, 58, 62, 65, 69, 74, 77, 82, 85, 87, 91, 93, 95, 111, 115, 119, 123, 133, 143, 155, 185, 203, 221, 287, 299, 319, 391, 407, 481, 933, 1555, 11507}
\solitem{3}{6, 10, 14, 15, 21, 22, 26, 33, 34, 35, 38, 39, 46, 51, 55, 57, 58, 62, 65, 69, 74, 77, 82, 85, 87, 91, 93, 95, 111, 115, 119, 123, 133, 145, 155, 161, 187, 221, 253, 259, 287, 407, 629, 667, 1509, 6539, 14587}
\solitem{4}{6, 10, 14, 15, 21, 22, 26, 33, 34, 35, 38, 39, 46, 51, 55, 57, 58, 62, 65, 69, 74, 77, 82, 85, 87, 91, 93, 95, 115, 119, 123, 129, 133, 143, 145, 155, 187, 213, 221, 287, 301, 559, 629, 781, 1247, 1591, 1633}
\solitem{5}{6, 10, 14, 15, 21, 22, 26, 33, 34, 35, 38, 39, 46, 51, 55, 57, 58, 62, 65, 69, 74, 77, 82, 86, 87, 91, 93, 95, 106, 111, 115, 123, 133, 155, 159, 161, 185, 203, 215, 253, 265, 583, 667, 731, 109591, 139823, 154939}
\solitem{6}{6, 10, 14, 15, 21, 22, 26, 33, 34, 35, 38, 39, 46, 51, 55, 57, 58, 62, 65, 69, 74, 77, 82, 86, 91, 93, 95, 111, 119, 122, 123, 133, 143, 145, 155, 161, 183, 187, 259, 287, 299, 319, 473, 481, 559, 671, 1591}
\solitem{7}{6, 10, 14, 15, 21, 22, 26, 33, 34, 35, 38, 39, 46, 51, 55, 57, 58, 62, 65, 69, 74, 77, 85, 86, 87, 91, 93, 94, 95, 111, 115, 119, 133, 141, 142, 145, 155, 187, 213, 517, 629, 799, 1247, 1591, 1633, 2021, 2627}
\solitem{8}{6, 10, 14, 15, 21, 22, 26, 33, 34, 35, 38, 39, 46, 51, 55, 57, 58, 62, 65, 69, 74, 77, 86, 87, 91, 93, 94, 95, 111, 115, 118, 119, 129, 133, 145, 155, 161, 185, 187, 329, 517, 667, 851, 1073, 1357, 1363, 2537}
\solitem{9}{6, 10, 14, 15, 21, 22, 26, 33, 34, 35, 38, 39, 46, 51, 55, 57, 58, 62, 65, 69, 74, 82, 85, 86, 87, 91, 93, 95, 106, 111, 123, 133, 145, 155, 159, 185, 203, 215, 253, 265, 287, 319, 493, 583, 731, 851, 1073}
\solitem{10}{6, 10, 14, 15, 21, 22, 26, 33, 34, 35, 38, 39, 46, 51, 55, 57, 58, 62, 65, 69, 77, 82, 85, 86, 87, 91, 93, 94, 95, 111, 115, 119, 133, 141, 142, 155, 187, 205, 213, 235, 493, 517, 1363, 1633, 1763, 1927, 2627}
\solitem{11}{6, 10, 14, 15, 21, 22, 26, 33, 34, 35, 38, 39, 46, 51, 55, 57, 58, 62, 65, 69, 77, 82, 85, 86, 87, 91, 93, 94, 95, 115, 118, 119, 123, 129, 133, 141, 155, 203, 287, 319, 391, 799, 1345, 1357, 2537, 2959, 12643}
\solitem{12}{6, 10, 14, 15, 21, 22, 26, 33, 34, 35, 38, 39, 46, 51, 55, 57, 58, 62, 65, 69, 77, 82, 85, 86, 87, 91, 93, 94, 95, 115, 118, 119, 123, 133, 143, 145, 155, 203, 221, 287, 391, 559, 1357, 2537, 12209, 18103, 19787}
\solitem{13}{6, 10, 14, 15, 21, 22, 26, 33, 34, 35, 38, 39, 46, 51, 55, 57, 58, 62, 65, 69, 77, 82, 85, 86, 87, 91, 93, 94, 95, 115, 119, 123, 129, 133, 141, 142, 145, 155, 187, 517, 799, 1207, 1247, 1633, 1763, 2021, 2911}
\solitem{14}{6, 10, 14, 15, 21, 22, 26, 33, 34, 35, 38, 39, 46, 51, 55, 57, 58, 62, 65, 69, 77, 82, 85, 86, 87, 91, 93, 94, 95, 118, 119, 123, 129, 133, 155, 161, 187, 213, 287, 329, 355, 493, 497, 517, 1357, 1363, 2537}
\solitem{15}{6, 10, 14, 15, 21, 22, 26, 33, 34, 35, 38, 39, 46, 51, 55, 57, 58, 62, 65, 69, 77, 82, 85, 86, 87, 91, 93, 94, 95, 118, 119, 123, 133, 145, 155, 161, 187, 203, 215, 287, 329, 493, 517, 1247, 1357, 1363, 2537}
\solitem{16}{6, 10, 14, 15, 21, 22, 26, 33, 34, 35, 38, 39, 46, 51, 55, 57, 58, 62, 65, 69, 77, 82, 85, 86, 87, 93, 94, 95, 115, 118, 119, 123, 129, 133, 141, 143, 145, 155, 287, 377, 391, 551, 799, 893, 1357, 1363, 2537}
\solitem{17}{6, 10, 14, 15, 21, 22, 26, 33, 34, 35, 38, 39, 46, 51, 55, 57, 58, 62, 65, 69, 77, 82, 85, 86, 87, 93, 94, 95, 115, 118, 119, 123, 129, 133, 143, 145, 155, 177, 287, 299, 391, 551, 611, 799, 893, 1711, 2537}
\solitem{18}{6, 10, 14, 15, 21, 22, 26, 33, 34, 35, 38, 39, 46, 51, 55, 57, 58, 62, 65, 69, 77, 82, 85, 86, 91, 93, 94, 95, 115, 119, 123, 133, 141, 142, 143, 155, 203, 221, 235, 287, 299, 355, 377, 391, 559, 2021, 2059}
\solitem{19}{6, 10, 14, 15, 21, 22, 26, 33, 34, 35, 38, 39, 46, 51, 55, 57, 58, 62, 65, 69, 77, 82, 85, 87, 91, 93, 94, 95, 111, 115, 118, 119, 123, 133, 142, 155, 187, 213, 287, 295, 329, 413, 493, 517, 1363, 1633, 2627}
\solitem{20}{6, 10, 14, 15, 21, 22, 26, 33, 34, 35, 38, 39, 46, 51, 55, 57, 58, 62, 65, 69, 77, 82, 86, 87, 91, 93, 94, 95, 118, 119, 123, 129, 133, 145, 155, 161, 177, 187, 203, 287, 299, 329, 377, 517, 1363, 1711, 2537}
\solitem{21}{6, 10, 14, 15, 21, 22, 26, 33, 34, 35, 38, 39, 46, 51, 55, 57, 58, 62, 65, 69, 77, 82, 86, 91, 93, 94, 95, 111, 115, 119, 123, 133, 141, 142, 143, 145, 155, 203, 213, 221, 235, 287, 493, 559, 1633, 2021, 2627}
\solitem{22}{6, 10, 14, 15, 21, 22, 26, 33, 34, 35, 38, 39, 46, 51, 55, 57, 58, 62, 65, 69, 77, 82, 87, 91, 93, 95, 106, 111, 115, 119, 122, 123, 133, 142, 155, 159, 183, 187, 203, 213, 265, 287, 319, 583, 671, 1633, 2627}
\solitem{23}{6, 10, 14, 15, 21, 22, 26, 33, 34, 35, 38, 39, 46, 51, 55, 57, 58, 62, 65, 74, 77, 82, 85, 87, 91, 93, 95, 111, 115, 119, 123, 129, 133, 143, 145, 155, 161, 221, 253, 259, 287, 299, 391, 473, 481, 1247, 1591}


% =====================================================================
\section{The computational hypotheses in Lean}\label{app:lean}
% =====================================================================

The exhaustiveness theorem of the search and the hypotheses of \cref{sec:lean}, verbatim from \file{lean/SemiprimeEgypt/SearchReal.lean}, \file{MainSearch.lean} and \file{Main.lean}.
The primed hypotheses \lean{Search46'} and \lean{Search47'} are obtained from \lean{Search46} and \lean{Search47} by restricting every quantifier over cores \code{C} to the cores in \code{searchCores K}, the output of the verified search; \lean{Search46'} is shown below, and \lean{Search47'} is \lean{Search47} with \code{∀ C ∈ searchCores 47,} in place of \code{∀ C : Finset ℕ,} in each of its four parts.
\begin{itemize}
\item \code{CoreSet C} says that every element of $C$ lies in $\Pset_\Sset$.
\item \code{Admissible K C} is \cref{def:admissible}.
\item \code{ShapeOK K C d b} is \cref{def:shape}.
\item \code{StarCand m C d P A} is \cref{def:candidate}.
\item \lean{mult}, \lean{starEdges} and \lean{starVal} are the multiplicities, elements and value of a star family.
\item \lean{compN} and \lean{compP\(_1\)}, \lean{compP\(_2\)} are $N$, $P_1$ and $P_2$ of \cref{prop:component}.
\item \lean{Sol23} is the family of \cref{app:solutions}.
\end{itemize}
\fontsize{8.3}{9.6}\selectfont
\begin{verbatim}
theorem core_mem_search {K : ℕ} (hK : K ≤ 48) {C : Finset ℕ} (hC : CoreSet C)
    (hA : Admissible K C) : edgesOf C ∈ (inst73 K).search (inst73 K).mkTables

def searchCores (K : ℕ) : List (Finset ℕ) :=
  ((inst73 K).search (inst73 K).mkTables).map coreOf

theorem core_mem_searchCores {T : Finset ℕ} (hT : ∀ n ∈ T, IsSP n) (hne : T.Nonempty)
    {K : ℕ} (hK : K ≤ 48) (hc : T.card ≤ K) (hz : ∃ z : ℤ, recipSum T = z) :
    core 73 T ∈ searchCores K

def Search46' : Prop :=
  ∀ C ∈ searchCores 46, CoreSet C → Admissible 46 C →
    (Uset 73 C).Nonempty ∧
      ∀ (d : ℕ → ℕ) (b : ℕ), ShapeOK 46 C d b → b = 0 ∧ ∀ P A, ¬ StarCand 10 C d P A

def Search46 : Prop :=
  ∀ C : Finset ℕ, CoreSet C → Admissible 46 C →
    (Uset 73 C).Nonempty ∧
      ∀ (d : ℕ → ℕ) (b : ℕ), ShapeOK 46 C d b → b = 0 ∧ ∀ P A, ¬ StarCand 10 C d P A

def Search47 : Prop :=
  (∀ C : Finset ℕ, CoreSet C → Admissible 47 C → Uset 73 C = ∅ → C.Nonempty → C ∈ Sol23) ∧
  (∀ C : Finset ℕ, CoreSet C → Admissible 47 C → (Uset 73 C).Nonempty →
      ∀ (d : ℕ → ℕ) (b : ℕ), ShapeOK 47 C d b → b ≤ 1) ∧
  (∀ (C : Finset ℕ) (d : ℕ → ℕ) (Ps : Finset ℕ) (A : ℕ → Finset ℕ),
      CoreSet C → Admissible 47 C → (Uset 73 C).Nonempty → ShapeOK 47 C d 0 →
      (∀ P ∈ Ps, StarCand 11 C d P (A P)) →
      (∀ q ∈ smallPrimes 73, mult Ps A q = d q) →
      recipSum C + starVal Ps A = 1 →
      C ∪ starEdges Ps A ∈ Sol23) ∧
  (∀ (C : Finset ℕ) (d : ℕ → ℕ) (Ps : Finset ℕ) (A : ℕ → Finset ℕ) (A₁ A₂ : Finset ℕ)
      (t X : ℕ),
      CoreSet C → Admissible 47 C → (Uset 73 C).Nonempty → ShapeOK 47 C d 1 →
      (∀ P ∈ Ps, StarCand 11 C d P (A P)) →
      A₁.Nonempty → A₂.Nonempty → A₁ ⊆ smallPrimes 73 → A₂ ⊆ smallPrimes 73 →
      (∀ q ∈ smallPrimes 73,
        mult Ps A q + (if q ∈ A₁ then 1 else 0) + (if q ∈ A₂ then 1 else 0) = d q) →
      recipSum C + starVal Ps A < 1 →
      0 < t →
      (t : ℚ) = (1 - recipSum C - starVal Ps A) * (∏ q ∈ A₁, (q : ℚ)) * ∏ q ∈ A₂, (q : ℚ) →
      0 < X → X ∣ compN A₁ A₂ t → t ∣ X + nA A₁ * ∏ q ∈ A₂, q →
      t ∣ compN A₁ A₂ t / X + nA A₂ * ∏ q ∈ A₁, q →
      compP₁ A₁ A₂ t X ≠ compP₂ A₁ A₂ t X →
      73 < compP₁ A₁ A₂ t X → 73 < compP₂ A₁ A₂ t X →
      (compP₁ A₁ A₂ t X).Prime → (compP₂ A₁ A₂ t X).Prime →
      compP₁ A₁ A₂ t X ∉ Ps → compP₂ A₁ A₂ t X ∉ Ps →
      C ∪ starEdges Ps A ∪ A₁.image (· * compP₁ A₁ A₂ t X) ∪
          A₂.image (· * compP₂ A₁ A₂ t X) ∪ {compP₁ A₁ A₂ t X * compP₂ A₁ A₂ t X} ∈ Sol23)
\end{verbatim}
\normalsize

\end{document}
